% Options for packages loaded elsewhere
% Options for packages loaded elsewhere
\PassOptionsToPackage{unicode}{hyperref}
\PassOptionsToPackage{hyphens}{url}
\PassOptionsToPackage{dvipsnames,svgnames,x11names}{xcolor}
%
\documentclass[
  letterpaper,
  DIV=11,
  numbers=noendperiod]{scrartcl}
\usepackage{xcolor}
\usepackage{amsmath,amssymb}
\setcounter{secnumdepth}{-\maxdimen} % remove section numbering
\usepackage{iftex}
\ifPDFTeX
  \usepackage[T1]{fontenc}
  \usepackage[utf8]{inputenc}
  \usepackage{textcomp} % provide euro and other symbols
\else % if luatex or xetex
  \usepackage{unicode-math} % this also loads fontspec
  \defaultfontfeatures{Scale=MatchLowercase}
  \defaultfontfeatures[\rmfamily]{Ligatures=TeX,Scale=1}
\fi
\usepackage{lmodern}
\ifPDFTeX\else
  % xetex/luatex font selection
\fi
% Use upquote if available, for straight quotes in verbatim environments
\IfFileExists{upquote.sty}{\usepackage{upquote}}{}
\IfFileExists{microtype.sty}{% use microtype if available
  \usepackage[]{microtype}
  \UseMicrotypeSet[protrusion]{basicmath} % disable protrusion for tt fonts
}{}
\makeatletter
\@ifundefined{KOMAClassName}{% if non-KOMA class
  \IfFileExists{parskip.sty}{%
    \usepackage{parskip}
  }{% else
    \setlength{\parindent}{0pt}
    \setlength{\parskip}{6pt plus 2pt minus 1pt}}
}{% if KOMA class
  \KOMAoptions{parskip=half}}
\makeatother
% Make \paragraph and \subparagraph free-standing
\makeatletter
\ifx\paragraph\undefined\else
  \let\oldparagraph\paragraph
  \renewcommand{\paragraph}{
    \@ifstar
      \xxxParagraphStar
      \xxxParagraphNoStar
  }
  \newcommand{\xxxParagraphStar}[1]{\oldparagraph*{#1}\mbox{}}
  \newcommand{\xxxParagraphNoStar}[1]{\oldparagraph{#1}\mbox{}}
\fi
\ifx\subparagraph\undefined\else
  \let\oldsubparagraph\subparagraph
  \renewcommand{\subparagraph}{
    \@ifstar
      \xxxSubParagraphStar
      \xxxSubParagraphNoStar
  }
  \newcommand{\xxxSubParagraphStar}[1]{\oldsubparagraph*{#1}\mbox{}}
  \newcommand{\xxxSubParagraphNoStar}[1]{\oldsubparagraph{#1}\mbox{}}
\fi
\makeatother


\usepackage{longtable,booktabs,array}
\usepackage{calc} % for calculating minipage widths
% Correct order of tables after \paragraph or \subparagraph
\usepackage{etoolbox}
\makeatletter
\patchcmd\longtable{\par}{\if@noskipsec\mbox{}\fi\par}{}{}
\makeatother
% Allow footnotes in longtable head/foot
\IfFileExists{footnotehyper.sty}{\usepackage{footnotehyper}}{\usepackage{footnote}}
\makesavenoteenv{longtable}
\usepackage{graphicx}
\makeatletter
\newsavebox\pandoc@box
\newcommand*\pandocbounded[1]{% scales image to fit in text height/width
  \sbox\pandoc@box{#1}%
  \Gscale@div\@tempa{\textheight}{\dimexpr\ht\pandoc@box+\dp\pandoc@box\relax}%
  \Gscale@div\@tempb{\linewidth}{\wd\pandoc@box}%
  \ifdim\@tempb\p@<\@tempa\p@\let\@tempa\@tempb\fi% select the smaller of both
  \ifdim\@tempa\p@<\p@\scalebox{\@tempa}{\usebox\pandoc@box}%
  \else\usebox{\pandoc@box}%
  \fi%
}
% Set default figure placement to htbp
\def\fps@figure{htbp}
\makeatother


% definitions for citeproc citations
\NewDocumentCommand\citeproctext{}{}
\NewDocumentCommand\citeproc{mm}{%
  \begingroup\def\citeproctext{#2}\cite{#1}\endgroup}
\makeatletter
 % allow citations to break across lines
 \let\@cite@ofmt\@firstofone
 % avoid brackets around text for \cite:
 \def\@biblabel#1{}
 \def\@cite#1#2{{#1\if@tempswa , #2\fi}}
\makeatother
\newlength{\cslhangindent}
\setlength{\cslhangindent}{1.5em}
\newlength{\csllabelwidth}
\setlength{\csllabelwidth}{3em}
\newenvironment{CSLReferences}[2] % #1 hanging-indent, #2 entry-spacing
 {\begin{list}{}{%
  \setlength{\itemindent}{0pt}
  \setlength{\leftmargin}{0pt}
  \setlength{\parsep}{0pt}
  % turn on hanging indent if param 1 is 1
  \ifodd #1
   \setlength{\leftmargin}{\cslhangindent}
   \setlength{\itemindent}{-1\cslhangindent}
  \fi
  % set entry spacing
  \setlength{\itemsep}{#2\baselineskip}}}
 {\end{list}}
\usepackage{calc}
\newcommand{\CSLBlock}[1]{\hfill\break\parbox[t]{\linewidth}{\strut\ignorespaces#1\strut}}
\newcommand{\CSLLeftMargin}[1]{\parbox[t]{\csllabelwidth}{\strut#1\strut}}
\newcommand{\CSLRightInline}[1]{\parbox[t]{\linewidth - \csllabelwidth}{\strut#1\strut}}
\newcommand{\CSLIndent}[1]{\hspace{\cslhangindent}#1}



\setlength{\emergencystretch}{3em} % prevent overfull lines

\providecommand{\tightlist}{%
  \setlength{\itemsep}{0pt}\setlength{\parskip}{0pt}}



 


\KOMAoption{captions}{tableheading}
\makeatletter
\@ifpackageloaded{caption}{}{\usepackage{caption}}
\AtBeginDocument{%
\ifdefined\contentsname
  \renewcommand*\contentsname{Table of contents}
\else
  \newcommand\contentsname{Table of contents}
\fi
\ifdefined\listfigurename
  \renewcommand*\listfigurename{List of Figures}
\else
  \newcommand\listfigurename{List of Figures}
\fi
\ifdefined\listtablename
  \renewcommand*\listtablename{List of Tables}
\else
  \newcommand\listtablename{List of Tables}
\fi
\ifdefined\figurename
  \renewcommand*\figurename{Figure}
\else
  \newcommand\figurename{Figure}
\fi
\ifdefined\tablename
  \renewcommand*\tablename{Table}
\else
  \newcommand\tablename{Table}
\fi
}
\@ifpackageloaded{float}{}{\usepackage{float}}
\floatstyle{ruled}
\@ifundefined{c@chapter}{\newfloat{codelisting}{h}{lop}}{\newfloat{codelisting}{h}{lop}[chapter]}
\floatname{codelisting}{Listing}
\newcommand*\listoflistings{\listof{codelisting}{List of Listings}}
\makeatother
\makeatletter
\makeatother
\makeatletter
\@ifpackageloaded{caption}{}{\usepackage{caption}}
\@ifpackageloaded{subcaption}{}{\usepackage{subcaption}}
\makeatother
\usepackage{bookmark}
\IfFileExists{xurl.sty}{\usepackage{xurl}}{} % add URL line breaks if available
\urlstyle{same}
\hypersetup{
  pdftitle={Cortext: A Three-Knob Adaptive Memory Architecture},
  pdfauthor={Gabriel Willen; Cortext Team},
  colorlinks=true,
  linkcolor={blue},
  filecolor={Maroon},
  citecolor={Blue},
  urlcolor={Blue},
  pdfcreator={LaTeX via pandoc}}


\title{Cortext: A Three-Knob Adaptive Memory Architecture}
\usepackage{etoolbox}
\makeatletter
\providecommand{\subtitle}[1]{% add subtitle to \maketitle
  \apptocmd{\@title}{\par {\large #1 \par}}{}{}
}
\makeatother
\subtitle{for Continuous Cognitive Processing}
\author{Gabriel Willen \and Cortext Team}
\date{2025-12-01}
\begin{document}
\maketitle


\section*{Abstract}\label{abstract}
\addcontentsline{toc}{section}{Abstract}

We present Cortext, a biologically-inspired adaptive memory system
governed by three continuous control parameters: Focus (F), Sensitivity
(S), and Stability (T). Unlike traditional memory architectures that
employ discrete operational modes, Cortext achieves developmental
progression through parameter-derived rate modulation, allowing behavior
to emerge continuously from the interaction of knob settings and
experiential mass. The architecture integrates established cognitive
science principles---including Cowan's working memory constraints and
Nader's reconsolidation dynamics---into a unified computational
framework. We derive most system parameters from the three primary knobs
through principled mathematical transformations, while explicitly
labeling the small set of fixed invariants (e.g., controller gains),
reducing reliance on hard-coded constants. The system demonstrates
self-calibrating priors that blend with evidence using
uncertainty-weighted Bayesian averaging, homeostatic threshold control
with effective sample size estimation, and graph-augmented retrieval
combining embedding similarity with semantic extraction. Experimental
analysis indicates the architecture maintains stable operation across
developmental phases while adapting write rates, decay dynamics, and
retrieval precision to environmental demands. This work contributes a
formally specified cognitive memory model suitable for implementation in
streaming AI systems requiring persistent, context-aware memory.

\textbf{Keywords:} cognitive architecture, adaptive memory, working
memory, episodic memory, semantic memory, knowledge graphs, homeostatic
control

\section{Introduction}\label{introduction}

Memory systems in artificial intelligence face a fundamental tension
between plasticity and stability. Systems that learn rapidly risk
catastrophic interference, while those that maintain stable
representations may fail to capture novel patterns (McCloskey and Cohen
1989). Biological memory systems resolve this tension through
sophisticated regulatory mechanisms that modulate learning rates, decay
dynamics, and retrieval thresholds in response to environmental demands
and internal state (McClelland, McNaughton, and O'Reilly 1995).

This paper introduces Cortext, a cognitive memory architecture that
addresses this stability-plasticity dilemma through three continuous
control parameters that govern all system behavior. Rather than
implementing discrete operational modes or hard-coded phase transitions,
Cortext achieves developmental progression through the continuous
interaction of parameter settings with accumulated experience. The
architecture draws on established findings from cognitive psychology and
neuroscience, including working memory capacity limits (Cowan 2001),
memory reconsolidation (Nader, Schafe, and Le Doux 2000), serial
position effects (Murdock Jr 1962), and emotional modulation of memory
(McGaugh 2004).

The core contribution of this work is a formally specified memory
architecture in which:

\begin{enumerate}
\def\labelenumi{\arabic{enumi}.}
\tightlist
\item
  Most tuneable parameters derive from three primary knobs (Focus,
  Sensitivity, Stability) through explicit mathematical transformations,
  while fixed invariants (e.g., controller gains) are labeled and
  justified.
\item
  System priors self-calibrate through uncertainty-weighted Bayesian
  blending with observed evidence.
\item
  Developmental phases emerge from annealed safety bounds and
  experiential mass accumulation, not explicit mode switching.
\item
  A knowledge graph layer enables semantic consolidation and
  graph-augmented retrieval.
\end{enumerate}

The remainder of this paper is organized as follows. Section 2 reviews
relevant literature. Section 3 presents the mathematical foundations
including notation, helper functions, and knob-derived parameters.
Section 4 details the core algorithms for Focus, Sensitivity, and
Stability adaptation. Section 5 describes structural metrics and
composite scoring. Section 6 covers dynamic thresholding and homeostatic
control. Section 7 presents the reinforcement and decay dynamics.
Section 8 describes advanced cognitive processes including working
memory, metacognition, and emotional consolidation. Section 9 details
the consolidation and graph integration system. Section 10 presents the
interrupt gate for streaming integration. Section 11 reports preliminary
experimental results. Section 12 discusses implementation considerations
and computational complexity. Section 13 concludes with limitations and
future directions.

\section{Related Work}\label{related-work}

\subsection{Working Memory Models}\label{working-memory-models}

The working memory component of Cortext draws primarily on Cowan (2001)
embedded-processes model, which posits a capacity limit of approximately
4±1 chunks for the focus of attention. This contrasts with Miller (1956)
earlier estimate of 7±2 items, which subsequent research has shown
conflates chunking with raw capacity (Cowan 2010). Our implementation
respects these empirically-derived constraints while allowing for
focus-dependent modulation within bounded ranges.

Baddeley (2000) multicomponent model informs our treatment of
maintenance and rehearsal processes, though we adopt a more unified
representational substrate based on distributed embeddings rather than
separate phonological and visuospatial stores. The episodic buffer
concept (Baddeley 2000) aligns with our approach to binding information
across modalities through shared vector spaces.

\subsection{Memory Consolidation}\label{memory-consolidation}

The consolidation mechanisms in Cortext reflect findings from the memory
reconsolidation literature (Nader, Schafe, and Le Doux 2000; Nader
2003). Reconsolidation theory posits that retrieved memories enter a
labile state during which they can be modified before restabilization.
Our architecture implements this through time-bounded lability windows
governed by the Stability parameter, with reconsolidation gain modulated
by both Sensitivity and contextual relevance.

The distinction between episodic and semantic memory (Tulving 1972)
motivates our two-tier storage approach: a streaming episodic buffer for
immediate experiences and a consolidated semantic graph for abstracted
knowledge. The consolidation process transforms high-redundancy episodic
clusters into summary nodes linked by typed semantic relations,
consistent with complementary learning systems theory (McClelland,
McNaughton, and O'Reilly 1995).

\subsection{Emotional Influences on
Memory}\label{emotional-influences-on-memory}

McGaugh (2004) extensive work on emotional modulation of memory
consolidation informs our treatment of affect-gated encoding. The
architecture implements emotional intensity as a threshold modifier,
consistent with findings that arousal enhances memory through
amygdala-mediated modulation of hippocampal encoding (LaBar and Cabeza
2006). We adopt a dimensional model of emotion (Russell 1980) with
valence and arousal as primary axes, projected from categorical emotion
embeddings.

\subsection{Adaptive Control Systems}\label{adaptive-control-systems}

The homeostatic threshold controller draws on classical control theory,
specifically proportional-integral approaches to setpoint maintenance
(Åström and Murray 2008). The use of exponentially-weighted moving
averages for rate estimation follows standard practice in adaptive
systems, while our effective sample size calculation for reliability
estimation extends techniques from sequential Monte Carlo methods (Liu
and Chen 1998).

\section{Mathematical Foundations}\label{sec-math}

\subsection{Notation and Primitives}\label{notation-and-primitives}

We establish the following notation used throughout this paper. Let ε(F,
S, T) = 10\^{}(−8 + 2T) denote a small stability-dependent constant for
numerical stability (we write ε as shorthand for ε(F, S, T)). All knob
values F, S, T lie in the closed interval {[}0, 1{]}.

Core mathematical primitives:

\begin{verbatim}
lerp(a, b, x) = a + (b − a) × x
clamp(v, lo, hi) = max(lo, min(v, hi))
sigmoid(z) = 1 / (1 + exp(−z))
EWMA(prev, x, α) = (1 − α) × prev + α × x
\end{verbatim}

Weight normalization and blending for combining multiple signals:

\begin{verbatim}
normalize(w) = w / max(sum(w), ε) # w is a weight vector
\end{verbatim}

Edge case: if sum(w) \textless{} ε, return uniform weights
(1/\textbar w\textbar) to avoid division by zero or invalid probability
distributions.

\begin{verbatim}
blend(values, weights) = Σᵢ values[i] × weights[i]
\end{verbatim}

Note: normalize() operates on weight vectors, not scalars. When used
with scalar expressions like lerp(), collect the scalars into a vector
first: normalize({[}lerp(\ldots), lerp(\ldots), \ldots{]}). The blend()
function assumes pre-normalized weights summing to 1.

For vectors, we define cosine similarity as cos(u, v) = u·v / (‖u‖ ×
‖v‖), and safe L2 normalization as l2\_normalize(v) = v / max(‖v‖, ε).
Shannon entropy is computed in nats: H(p) = −Σᵢ pᵢ ln(pᵢ).

The temporal decay function follows exponential dynamics with
configurable half-life:

\begin{verbatim}
decay(x, τ_half, Δt) = x × exp(−ln(2) × Δt / max(τ_half, τ_min(T)))
\end{verbatim}

where τ\_min(T) = lerp(60, 300, T) seconds provides a stability-derived
floor to prevent numerical instability from near-zero half-lives.

\subsubsection{Units Convention}\label{units-convention}

To ensure consistency and avoid unit mismatch errors, the following
conventions apply throughout:

\textbf{Input timestamps:} All input timestamps (t) are specified in
milliseconds since epoch, consistent with standard system time APIs.

\textbf{Internal time representation:} All internal time calculations
operate in seconds. We define:

\begin{verbatim}
now_s() = system_time_ms / 1000  # returns seconds
now_ms() = system_time_ms         # returns milliseconds
to_s(ts_ms) = ts_ms / 1000       # ms → seconds conversion
\end{verbatim}

\textbf{Stored timestamps:} All stored timestamps are milliseconds since
Unix epoch.

\textbf{Naming contract (canonical):} Use the variable names below
consistently throughout this document. Units: stored timestamps are
integers in milliseconds (commonly suffixed *\_ts, and also appearing as
timestamp/created\_at/last\_rate\_timestamp); derived time intervals in
seconds use the \emph{\emph{s suffix. Accumulator variables: t\_start,
last\_signal\_ts, last\_write\_ts, drift\_acc, eta\_acc,
coherence\_prev, emo\_max, arousal\_sum, drift\_accum,
drift\_at\_last\_interrupt, drift\_acc\_pacing, x\_last\_check. Global
variables: signals\_processed, u\_uncertainty, mood\_vector,
last\_mood\_ts, theta\_dynamic, theta\_target, hysteresis, m\_rate,
rho\_hat\_prev, dt\_ema, rate\_ticks, reliability,
last\_rate\_timestamp, last\_retrieval\_ts, retention\_ema,
last\_embedding, x\_pred\_ema. Weight naming rule: weight}} and
\emph{\emph{weight variables (e.g., weight\_relevance, mismatch\_weight,
weight\_surprise) are control parameters; w}} variables (e.g.,
w\_relevance, w\_mismatch, \ldots{} w\_arousal) are composite-score
blender weights.

\textbf{Drift naming note:} \texttt{drift\_acc} tracks
within-accumulator centroid drift (context drift used for boundary
decisions); \texttt{drift\_accum} tracks cumulative per-signal embedding
drift used for interrupt refractory pressure. Treat these as distinct
signals.

\textbf{Time-index convention:} Per-step computed values use a
\texttt{\_t} suffix (e.g., score\_t, theta\_dynamic\_t). Retained state
uses the bare name (e.g., theta\_dynamic). When a step proposes a new
value, compute the \texttt{\_t} variant and assign the bare name at end
of step (see Appendix D).

\textbf{Time deltas:} All Δt values, elapsed times (mem\_elapsed,
signal\_gap, idle\_for), and time comparisons operate in seconds:

\begin{verbatim}
Δt ← now_s() − to_s(last_rate_timestamp)   # seconds
mem_elapsed ← now_s() − to_s(t_start)     # seconds
signal_gap ← now_s() − to_s(last_signal_ts)  # seconds
\end{verbatim}

\textbf{Time constants:} All time-related constants are specified with
explicit units (e.g., τ\_min(T) in seconds). Threshold limits like
max\_mem\_time(T) and gap\_scale(T) return values in seconds (gap\_scale
multiplies dt\_ema to form gap\_ref\_s).

\subsubsection{No Fixed Behavioral
Constants}\label{no-fixed-behavioral-constants}

All behavioral thresholds, gains, and caps are derived from the three
knobs (F, S, T) and/or state. Fixed numeric values are permitted only
for numerical stability floors (e.g., ε(F, S, T)) and unit conversions.
This keeps the system's qualitative behavior entirely governed by the
three knobs, consistent with the human-memory objective.

\subsubsection{Buffers}\label{sec-buffers}

The specification uses distinct streaming buffers:

\textbf{signal\_stream:} The raw per-signal embedding stream x\_t used
for scoring, uncertainty estimation, threshold/rate updates, and
accumulator updates. Any context slices operate on signal\_stream.

\textbf{score\_stream:} The per-signal scalar score stream. At each
step, append score\_t to score\_stream. Any score lookbacks operate on
score\_stream.

\textbf{memory\_stream:} The stream of written memory representatives
(e.g., e\_rep/μ\_acc for completed units) used for retrieval and
interrupt-gate context. Sections that reference
recent\_memory\_centroids operate on memory\_stream.

\textbf{recent\_memory\_centroids:} A bounded deque of recent memory
representatives used by the interrupt gate.

\begin{verbatim}
win_mem_ctx(T) = round(lerp(4, 32, T))  # max memories in interrupt context
On successful write_memory: recent_memory_centroids.append(e_rep); recent_memory_centroids ← tail(recent_memory_centroids, win_mem_ctx(T))
\end{verbatim}

\subsection{The Three-Knob Philosophy}\label{the-three-knob-philosophy}

Cortext is governed by three continuous control parameters, each
representing a distinct dimension of cognitive regulation:

\textbf{Focus (F ∈ {[}0, 1{]}):} Perceptual selectivity and precision.
Higher Focus narrows attention, increases relevance weighting, and
reduces retrieval breadth. Focus modulates the trade-off between
exploitation of known-relevant information and exploration of
potentially useful context.

\textbf{Sensitivity (S ∈ {[}0, 1{]}):} Plasticity and affective gain.
Higher Sensitivity accelerates learning, increases emotional and novelty
responsiveness, and raises write-rate targets. Sensitivity governs how
readily the system captures novel information and responds to salient
stimuli.

\textbf{Stability (T ∈ {[}0, 1{]}):} Temporal persistence and inertia.
Higher Stability lengthens memory half-lives, widens hysteresis bands,
slows adaptive updates, and tightens safety bounds over time. Stability
controls the resistance to change and the preservation of established
knowledge.

A central design principle is that knobs set rates rather than modes.
Behavioral differences emerge continuously from parameter interactions;
there are no hard-coded phase transitions or discrete operational
states.

\subsection{Knob-Derived Parameters}\label{knob-derived-parameters}

Most system tunables derive from the three primary knobs. This section
catalogs the key derivations; fixed invariants are explicitly labeled in
their respective sections.

\subsubsection{Context Windows and Temporal
Scales}\label{context-windows-and-temporal-scales}

\begin{verbatim}
n_ctx(T) = round(lerp(32, 256, T))
win_score(T) = round(lerp(20, 120, T))
win_rate_s(T) = round(lerp(60, 300, T))
\end{verbatim}

The context window n\_ctx determines how many recent items inform
relevance computation. The scoring window win\_score controls the
lookback for variance estimation and percentile calculation. The rate
window win\_rate\_s specifies the temporal horizon (in seconds) for
write-rate measurement.

\subsubsection{Half-Life and Decay}\label{half-life-and-decay}

Memory half-life follows a log-scale mapping to span multiple orders of
magnitude:

\begin{verbatim}
τ_min = 120.0 seconds (2 minutes)
τ_max = 43200.0 seconds (12 hours)
base_half_life(T) = exp(ln(τ_min) + T × ln(τ_max / τ_min))
\end{verbatim}

This exponential mapping ensures that low Stability yields half-lives
near 2 minutes while high Stability approaches 12 hours, with smooth
interpolation across the range.

\subsubsection{Hysteresis and Rate
Targets}\label{hysteresis-and-rate-targets}

\begin{verbatim}
band_min = 0.02; band_max = 0.25
base_band(T) = lerp(band_min, band_max, T)
r_min = 0.2; r_max = 5.0  (writes per minute)
base_rate(S) = lerp(r_min, r_max, S)
\end{verbatim}

The hysteresis band prevents oscillation in threshold-crossing
decisions. Write-rate targets establish homeostatic setpoints for the
threshold controller.

These ranges are canonical knob maps; no additional fixed thresholds are
introduced beyond the three knobs.

\subsubsection{Experiential Mass and
Maturity}\label{experiential-mass-and-maturity}

The system tracks accumulated experience through a maturity function
that governs the annealing of safety bounds:

\begin{verbatim}
τ_m(T) = lerp(10.0, 200.0, T)
maturity(t) = 1 − exp(−count / τ_m(T))
\end{verbatim}

where count is the total number of signals processed. This produces
asymptotic approach to unit maturity, with higher Stability slowing the
progression to reflect greater conservatism.

Safety bounds on the dynamic threshold anneal with maturity:

\begin{verbatim}
T_min(t) = lerp(0.01, 0.05, maturity(t))
T_max(t) = lerp(0.99, 0.95, maturity(t))
max_ΔT_per_min(t) = lerp(0.30, 0.10, maturity(t))
\end{verbatim}

Early operation permits wide threshold excursions; mature operation
constrains movement to a narrower band.

\subsection{Uncertainty Estimation}\label{uncertainty-estimation}

Uncertainty u(t) ∈ {[}0, 1{]} modulates learning rates and evidence
weighting. The raw uncertainty estimate blends multiple signals:

\begin{verbatim}
var_score_max(S) = lerp(0.15, 0.35, S)
recent_scores ← tail(score_stream, win_score(T))
if |recent_scores| < 2:
    var_recent_norm ← 0
else:
    var_recent_norm ← clamp(var(recent_scores) / var_score_max, 0, 1)
coherence_complement = 1 − coherence_struct_t  # uses structural coherence
\end{verbatim}

When prediction error signals are available, novelty and surprisal are
blended:

\begin{verbatim}
novelty_surprise = blend([novelty_t, surprisal_t],
                        weights = normalize([S, 1 − T]))
\end{verbatim}

The final raw uncertainty combines these components with knob-derived
weights:

Normative note (MUST): focus\_spread\_t is the per-step derived
focus-spread metric from Section~\ref{sec-focus-spread} and is available
each step when computing u(t).

\begin{verbatim}
weights_u = normalize([S, F, 1 − T, S × (1 − T)])
u_raw(t) = clamp(blend([var_recent_norm, focus_spread_t,
                        coherence_complement, novelty_surprise],
                       weights = weights_u), 0, 1)
\end{verbatim}

Smoothed uncertainty applies EWMA with a stability-dependent rate:

\begin{verbatim}
α_u(T) = 0.10 + (1 − T) × 0.60
u(t) = EWMA(u(t−1), u_raw(t), α = α_u(T))
\end{verbatim}

When structural metrics are unavailable, the fallback is u\_raw(t) = 1 −
maturity(t), ensuring high uncertainty during early operation.

\section{Core Adaptation Algorithms}\label{sec-core-adaptation}

This section presents the algorithms governing adaptation along each of
the three primary dimensions. Each algorithm consists of a prior
computation (executed at initialization) and a dynamic update (executed
per signal).

\subsection{Focus-Driven Selectivity}\label{focus-driven-selectivity}

Focus governs perceptual selectivity through relevance weighting and
attention width.

\subsubsection{Focus Priors}\label{focus-priors}

Given Focus knob F ∈ {[}0, 1{]}, initialize Focus control variables:

\begin{verbatim}
weight_relevance = sigmoid(2F − 1)
coverage_gain_floor = 0.3 + 0.7F
mismatch_weight = 1 − F
attention_width = lerp(π, 0.1π, F)
\end{verbatim}

The attention width (in radians) controls the angular spread of the
receptive field in embedding space. High Focus produces narrow attention
(0.1π), while low Focus permits broad capture (π).

\subsubsection{Dynamic Focus Update}\label{dynamic-focus-update}

At each signal event t with input embedding x\_t:

\begin{verbatim}
recent_context ← tail(signal_stream, n_ctx(T))
if |recent_context| == 0:
    μ_ctx ← 0_vector; observed_cosine ← 0  # map01(0)=0.5
else:
    μ_ctx ← mean(recent_context)
    observed_cosine ← cos(x_t, μ_ctx)
weight_relevance_t ← EWMA(weight_relevance,
                      map01(observed_cosine), α = α_F(t))
\end{verbatim}

\begin{verbatim}
weight_relevance ← weight_relevance_t
\end{verbatim}

where map01(z) = clamp((z + 1) / 2, 0, 1) transforms cosine values from
{[}−1, 1{]} to {[}0, 1{]}.

The learning rate α\_F(t) is modulated by uncertainty:

\begin{verbatim}
α_min_F = 0.05; α_span_F = 0.45
α_F(t) = α_min_F + F × α_span_F × u(t)
\end{verbatim}

High uncertainty increases learning rate, allowing faster adaptation
when the environment is volatile. The Focus knob scales the uncertainty
responsiveness.

\subsection{Sensitivity-Driven
Plasticity}\label{sensitivity-driven-plasticity}

Sensitivity governs learning speed, emotional responsiveness, and
novelty capture.

\subsubsection{Sensitivity Priors}\label{sensitivity-priors}

Given Sensitivity knob S ∈ {[}0, 1{]}, initialize Sensitivity control
variables:

\begin{verbatim}
rate_target = base_rate(S)  # writes/min
weight_novelty = 0.3 + 0.7S
weight_surprise = 0.2 + 0.8S
weight_valence = 0.4 + 0.6S
weight_arousal = S
emotion_gain = exp(1.5S)
score_gain = exp(2S)
\end{verbatim}

\subsubsection{Emotional Projection}\label{sec-emotional-projection}

When emotion category centroids are available, the system projects input
embeddings onto a discrete emotion space C = \{anger, fear, joy, love,
sadness, surprise\}. categories inspired by Russell (1980) circumplex
model are projected onto the valence-arousal plane:coordinates:

\begin{verbatim}
v_map = {anger: −0.9, fear: −0.8, sadness: −0.9,
         joy: +0.9, love: +0.8, surprise: 0.0}
a_map = {anger: +0.9, fear: +0.9, sadness: +0.3,
         joy: +0.6, love: +0.5, surprise: +0.8}
\end{verbatim}

The projection procedure:

\begin{verbatim}
raw_cos_c ← cos(x_t, centroids[c]) for each c ∈ C
if all raw_cos_c ≤ 0:
    p_c ← uniform(1/6)  # ensures downstream mood update is well-defined
    emotion_intensity_t ← 0; valence_t ← 0.5; arousal_t ← 0
else:
    logits_c ← max(0, raw_cos_c)
    β(S) = 4 + 8S  # softmax inverse temperature
    p_c ← softmax(β(S) × logits_c)
    peak ← max_c(p_c)
    confidence ← 1 − H(p_c) / ln(6)
    emotion_intensity_t ← sqrt(peak × confidence)
    valence_t ← (Σ_c p_c × v_map[c] + 0.9) / 1.8
    arousal_t ← clamp(Σ_c p_c × a_map[c], 0, 1)
\end{verbatim}

The emotion intensity combines peak probability with distributional
confidence via geometric mean, providing a measure that is high only
when a single emotion dominates with high certainty.

\subsubsection{Threshold Modulation from
Emotion}\label{sec-emotional-threshold}

Emotional activation loosens write thresholds to capture salient
moments:

\begin{verbatim}
κ_emo ← κ_base × S  # where κ_base = 0.10
ΔThreshold_emotion_t ← −κ_emo × emotion_intensity_t ×
                        (0.5 + 0.5 × arousal_t)
\end{verbatim}

The emotional state acts as a modulator for memory
encoding/consolidation, following McGaugh (2004). High arousal and
valence magnitude increase the likelihood of threshold crossings.

\subsubsection{Mood Integration}\label{sec-mood}

Distinct from instantaneous emotion, the mood state M\_t ∈ ℝ⁶ maintains
a persistent background affective tone as a 6-dimensional vector (one
component per emotion category from
Section~\ref{sec-emotional-projection}):

\begin{verbatim}
α_mood(S) = lerp(0.01, 0.20, S)  # reactivity
half_life_mood(T) = lerp(30, 600, T)  # seconds
Δt_mood ← now_s() − to_s(last_mood_ts)
λ_mood(Δt_mood, T) ← exp(−ln(2) × Δt_mood / max(half_life_mood(T), ε))
e_t ← p_c − (1/6)  # centered 6D vector (can be negative)
M_t = λ_mood(Δt_mood, T) × M_{t−1} + α_mood(S) × e_t
M_t ← clamp_elementwise(M_t, −1.0, 1.0)  # per-component
last_mood_ts ← now_ms()  # update timestamp after mood update (ms)
\end{verbatim}

Note: when all raw\_cos\_c ≤ 0, p\_c is uniform and e\_t = 0, so mood
only decays. Because e\_t is centered around zero, M\_t can have both
positive and negative components, reflecting sustained elevation or
suppression relative to baseline. The mood state provides a separate
threshold bias via its normalized magnitude:

\begin{verbatim}
κ_mood ← κ_base × S
m_norm ← ‖M_t‖ / √6  # max norm when all components at 1
ΔThreshold_mood_t ← −κ_mood × clamp(m_norm, 0, 1)
\end{verbatim}

\subsection{Stability-Driven
Persistence}\label{stability-driven-persistence}

Stability governs temporal dynamics through half-life, decay rates, and
hysteresis.

\subsubsection{Stability Priors}\label{stability-priors}

Given Stability knob T ∈ {[}0, 1{]}, initialize Stability control
variables:

\begin{verbatim}
hysteresis = lerp(0.02, 0.25, T)
half_life = base_half_life(T)
rate_decay = lerp(0.60, 0.98, T)
periphery_half_life = clamp(0.5 × half_life,
                                   τ_min, τ_max)
drift_weight = 0.5 × (1 − T)
\end{verbatim}

\subsubsection{Dynamic Stability Update}\label{sec-stability-update}

The stability update uses uncertainty-modulated learning rate and a
stability-derived retention window:

\begin{verbatim}
α_min_T = 0.02; α_span_T = 0.18
α_T(t) = α_min_T + (1 − T) × α_span_T × u(t)
win_ret(T) = round(lerp(10, 50, T))  # retention history window size
\end{verbatim}

At each signal event, compute retention statistics and adjust half-life:

\begin{verbatim}
active_memories ← {m | strength(m) ≥ periphery_cutoff(T)}
observed_retention ← mean_age(active_memories)
retention_ema_t ← EWMA(retention_ema_{t−1},
                       observed_retention, α = α_T(t))
\end{verbatim}

\begin{verbatim}
retention_ema ← retention_ema_t
\end{verbatim}

Definitions:

\begin{verbatim}
age_s(m) ← now_s() − to_s(m.created_at)  # fallback: use start_ts if created_at unset
mean_age(active_memories) ← mean_{m ∈ active_memories} age_s(m); if empty, return 0
retention_history ← append(observed_retention); keep last win_ret(T) values
\end{verbatim}

Compute z-score relative to recent retention history:

\begin{verbatim}
last_win_ret ← tail(retention_history, win_ret(T))
μ_ret ← mean(last_win_ret)
σ_ret ← max(std(last_win_ret), 1.0)
zscore_ret ← clamp((observed_retention − μ_ret) / σ_ret, −3, +3)
\end{verbatim}

The target half-life incorporates feedback adjustment from the stability
feedback mechanism (Section~\ref{sec-feedback}):

\begin{verbatim}
stability_adj ← ΔHalfLife_adj_t if provided else 0
target_half_life_t ← clamp(base_half_life(T) ×
                          (1 + 0.25 × zscore_ret + stability_adj),
                          τ_min, τ_max)
half_life_t ← EWMA(half_life_{t−1}, target_half_life_t,
                   α = α_T(t))
\end{verbatim}

\begin{verbatim}
half_life ← half_life_t
\end{verbatim}

\section{Structural Metrics and Composite
Scoring}\label{sec-structural-metrics}

Context definitions (used throughout Table 1 and structural metrics):

\begin{verbatim}
recent_context ← tail(signal_stream, n_ctx(T))
if |recent_context| == 0:
    μ_ctx ← 0_vector  # define cos(x_t, μ_ctx)=0 so map01(cos)=0.5
else:
    μ_ctx ← mean(recent_context)
\end{verbatim}

\subsection{Embedding-Derived Metrics}\label{embedding-derived-metrics}

\subsubsection{Structural Coherence}\label{sec-struct-coherence}

Structural coherence (coherence\_struct) measures integration of the
current signal with the broader context window. This metric is distinct
from memory coherence (coherence\_mem, defined in
Section~\ref{sec-drift-coherence}) which tracks within-memory
similarity:

\begin{verbatim}
coh_neutral(T) = lerp(0.45, 0.55, T)
if |recent_context| < 2:
    coherence_struct_t ← coh_neutral(T)  # neutral, stability-derived
else:
    raw ← var([cos(x_t, c) for c in recent_context])
    coherence_struct_t ← 1 − clamp(raw, 0, 1)  # range [0, 1]
\end{verbatim}

High structural coherence (low variance in similarities) indicates the
signal fits consistently with context. The effective Focus is modulated:
F\_eff = F × (0.5 + 0.5 × coherence\_struct\_t).

\textbf{Effective Focus usage:} F\_eff is a diagnostic modulation of
Focus. Unless a formula explicitly references F\_eff, the specification
uses raw F (including all Table 1 metric formulas and knob-derived
controls).

\subsubsection{Focus Spread}\label{sec-focus-spread}

Focus spread quantifies the entropy of attention over nearest neighbors:

\begin{verbatim}
k ← k_neighbors(T) = round(lerp(8, 32, T))
\end{verbatim}

Normative note (MUST): kNN\_similarities MUST be computed by querying
memory\_stream with q = x\_t and k = k\_neighbors(T) (not
recent\_context).

\begin{verbatim}
if |memory_stream| == 0:
    focus_spread_t ← 0  # cold-start fallback
else:
    k_eff ← min(k, |memory_stream|)
    if k_eff < 2:
        focus_spread_t ← 0  # avoid degenerate entropy
    else:
        kNN_similarities ← topK(vector_search(x_t, k=k_eff))
        # attention_width sharpens or flattens similarities
        kNN_similarities ← clamp(kNN_similarities × (π / attention_width), −1, 1)
        p ← softmax(kNN_similarities)
        focus_spread_t ← H(p) / ln(k_eff)
\end{verbatim}

Values near 1 indicate diffuse attention; values near 0 indicate
concentrated attention. The effective Focus is further modulated: F\_eff
← F\_eff × (1 − focus\_spread\_t).

\subsubsection{Trajectory Drift}\label{trajectory-drift}

Drift measures directional change in context centroids:

\begin{verbatim}
k_ctx(T) = round(lerp(3, 10, T))  # step lag for drift
ctx_t = signal_stream[t − n_ctx(T) + 1 : t]  # inclusive end
ctx_{t−k} = signal_stream[t − k_ctx(T) − n_ctx(T) + 1 : t − k_ctx(T)]
drift_vec_t ← l2_normalize(mean(ctx_t)) −
              l2_normalize(mean(ctx_{t−k}))
drift_mag_t ← ‖drift_vec_t‖
\end{verbatim}

Since both centroids are unit-normalized, drift\_mag\_t ∈ {[}0, 2{]}. A
threshold defines an informational drift-boundary signal:

\begin{verbatim}
drift_threshold ← lerp(0.10, 0.35, T)
drift_boundary_t ← (drift_mag_t > drift_threshold)
\end{verbatim}

Normative note (MUST): drift\_boundary\_t is informational and MUST NOT
trigger a memory flush on its own. Memory flush decisions are defined by
should\_flush in Section~\ref{sec-boundary}.

\subsubsection{Real-Time Prediction Error
(EMA)}\label{real-time-prediction-error-ema}

We model real-time expectation using an exponential moving average (EMA)
of the signal stream. The predictor state \texttt{x\_pred\_ema}
represents the current expectation; low Stability (T) adapts quickly,
high Stability adapts slowly.

We treat orthogonal or anti-correlated vectors (cos ≤ 0) as maximum
surprise.

\begin{verbatim}
β_pred(T) = lerp(0.25, 0.02, T)
err_ref(S) = lerp(0.25, 0.05, S)
k_surprise(S,T) = lerp(6, 14, S) × lerp(1.1, 0.9, T)

if x_pred_ema is unset:
    x_pred_ema ← x_t  # initialize expectation
    surprisal_t ← 0
else:
    prediction_error_t ← 1 − max(0, cos(x_t, x_pred_ema))
    surprisal_t ← sigmoid((prediction_error_t − err_ref(S)) * k_surprise(S,T))
    x_pred_ema ← l2_normalize((1 − β_pred(T)) × x_pred_ema +
                               β_pred(T) × x_t)
\end{verbatim}

This formulation captures prediction failure directly from the stream,
without requiring an explicit kinematic trend model.

\subsection{Composite Score
Computation}\label{composite-score-computation}

The system computes 12 metrics that blend into a composite write score:

\begin{longtable}[]{@{}
  >{\raggedright\arraybackslash}p{(\linewidth - 4\tabcolsep) * \real{0.3333}}
  >{\raggedright\arraybackslash}p{(\linewidth - 4\tabcolsep) * \real{0.3333}}
  >{\raggedright\arraybackslash}p{(\linewidth - 4\tabcolsep) * \real{0.3333}}@{}}
\toprule\noalign{}
\begin{minipage}[b]{\linewidth}\raggedright
Metric
\end{minipage} & \begin{minipage}[b]{\linewidth}\raggedright
Knob
\end{minipage} & \begin{minipage}[b]{\linewidth}\raggedright
Expression
\end{minipage} \\
\midrule\noalign{}
\endhead
\bottomrule\noalign{}
\endlastfoot
Relevance & ↑F &
\texttt{relevance\_t\ =\ clamp(map01(cos(x\_t,\ μ\_ctx))\ ×\ (0.5\ +\ 0.5F),\ 0,\ 1)} \\
Mismatch & ↓F, ↑S & \texttt{(1\ −\ F)\ ×\ S\ ×\ novelty\_t} \\
Surprise & ↑S, ↓T & \texttt{surprisal\_t\ ×\ S\ ×\ (1\ −\ 0.5T)} \\
Rarity & ↑F, ↓T &
\texttt{rarity\_t\ ×\ (0.5\ +\ 0.5F)\ ×\ (1\ −\ 0.2T)} \\
Drift & ↓T & \texttt{(drift\_mag\_t\ /\ 2)\ ×\ (1\ −\ T)} \\
Utility & ↑F, ↓S & \texttt{ΔSSE\ ×\ (0.5\ +\ 0.5F)\ ×\ (1\ −\ 0.3S)} \\
Salience & F, S &
\texttt{(rarity\_t\ +\ novelty\_t)\ /\ 2\ ×\ (F\ +\ S)\ /\ 2} \\
Valence & S, ↓T & \texttt{valence\_t} \\
Arousal & S, ↓T & \texttt{arousal\_t} \\
Contradiction & ↑S, ↓F & \texttt{max(0,\ S\ −\ F)} \\
Periphery & ↑T & \texttt{(1\ −\ relevance\_t)\ ×\ T} \\
Coverage & ↑F & \texttt{F\ ×\ relevance\_t} \\
\end{longtable}

\emph{Table 1: Metric definitions and knob dependencies. Arrows indicate
direction of influence.}

\subsection{Metric Weight Blending}\label{metric-weight-blending}

Metric weights adapt online using recursive least squares (RLS) to
minimize prediction error between composite scores and observed
outcomes. Initial weights derive from bootstrap coefficients:

Blender weights are maintained as variables w\_* (do not confuse these
with control weights like weight\_relevance or mismatch\_weight). We
denote the 12-element blender weight vector as:

\begin{verbatim}
W_blend = [w_relevance, w_mismatch, w_surprise, w_rarity, w_drift,
          w_utility, w_salience, w_valence, w_arousal, w_contradiction,
          w_periphery, w_coverage]
\end{verbatim}

The index i used below (e.g., w\_bootstrap{[}i{]}, w\_rls{[}i{]})
follows this ordering.

\begin{verbatim}
w_bootstrap[i] ← sigmoid(c_F[i]×F + c_S[i]×S + c_T[i]×T + d_i)
\end{verbatim}

Before normalization, control weights modulate the blender weights:

\begin{verbatim}
w_relevance ← w_relevance × weight_relevance
w_mismatch  ← w_mismatch  × mismatch_weight
w_surprise  ← w_surprise  × weight_surprise
w_valence   ← w_valence   × weight_valence × emotion_gain
w_arousal   ← w_arousal   × weight_arousal × emotion_gain
w_coverage  ← w_coverage  × coverage_gain_floor
\end{verbatim}

Composite score scaling applies after weight normalization:

\begin{verbatim}
score_t ← clamp(score_raw × score_gain, 0, 1)
\end{verbatim}

Bootstrap coefficient defaults (canonical; used for initialization):

\begin{longtable}[]{@{}lllll@{}}
\toprule\noalign{}
Metric & c\_F & c\_S & c\_T & d \\
\midrule\noalign{}
\endhead
\bottomrule\noalign{}
\endlastfoot
relevance & 1.4 & 0.0 & 0.4 & −1.0 \\
mismatch & −1.0 & 1.0 & 0.0 & −0.5 \\
surprise & 0.0 & 1.5 & −0.5 & −0.75 \\
rarity & 0.9 & 0.0 & −0.3 & 0.05 \\
drift & 0.0 & 0.0 & −1.0 & 0.0 \\
utility & 0.85 & −0.45 & 0.0 & 0.075 \\
salience & 1.0 & 1.0 & 0.0 & −1.0 \\
valence & 0.0 & 1.02 & −0.42 & −0.11 \\
arousal & 0.0 & 1.8 & −0.2 & −0.9 \\
contradiction & −2.0 & 2.0 & 0.0 & −1.0 \\
periphery & 0.0 & 0.0 & 1.0 & −1.0 \\
coverage & 2.0 & 0.0 & 0.0 & −1.0 \\
\end{longtable}

These coefficients define a canonical initialization map from the three
knobs; they do not introduce additional runtime thresholds and may be
refit from human data without changing the rest of the system.

\subsubsection{RLS Weight Adaptation
(Canonical)}\label{rls-weight-adaptation-canonical}

RLS fits the blender weights directly in \textbf{weight-space}, using
the current metrics as predictors and an observed outcome signal as the
target. State is retained as a weight vector \texttt{w} and covariance
matrix \texttt{P}.

Target and predictor:

\begin{verbatim}
x_t ← [metric_i] (each metric clamped to [0,1])
y_t ← relevance_t  # observed outcome (in [0,1])
\end{verbatim}

Update cadence and forgetting:

\begin{verbatim}
k_update = round(lerp(1, 8, T))  # update every k_update signals
λ(T) = 0.90 + 0.09T              # forgetting factor
\end{verbatim}

RLS update (performed when step \% k\_update == 0):

\begin{verbatim}
y_hat ← wᵀ x_t
e_t ← y_t − y_hat
K ← (P x_t) / (λ + x_tᵀ P x_t)
w ← clamp(w + K × e_t, 0, 1)
P ← (I − K x_tᵀ) P / λ
\end{verbatim}

Note: RLS fits \textbf{unnormalized} weights \texttt{w}; normalization
happens later in composite score computation.

Initialization:

\begin{verbatim}
w_0 ← w_bootstrap   # from the coefficient table above
P_init(T) = lerp(500, 2000, 1 − T)
P_0 ← diag(P_init(T))
\end{verbatim}

Interpretation-only window size:

\begin{verbatim}
N ← round(lerp(64, 512, T))  # effective fitting window (interpretation only)
\end{verbatim}

The fitted weights blend with bootstrap weights based on RLS confidence:

\begin{verbatim}
τ_rls ← lerp(20.0, 80.0, T)
t ← signals_processed  # global step count (maturity)
confidence_rls ← 1 − exp(−t / τ_rls)
w_rls01[i] ← clamp(w[i], 0, 1)  # constrain fitted weights to mixture range
weight_i(t) ← clamp((1 − confidence_rls) × w_bootstrap[i] +
                      confidence_rls × w_rls01[i], 0, 1)
\end{verbatim}

\subsubsection{Score Normalization}\label{score-normalization}

Composite score computation requires careful normalization:

\begin{verbatim}
for each metric i:
    m01[i] = clamp(metric[i], 0, 1)
weight_sum ← Σ weights[i]
if weight_sum < ε: return 0
weights_norm[i] ← weights[i] / weight_sum
score ← clamp(Σ weights_norm[i] × m01[i], 0, 1)
\end{verbatim}

Invariant (MUST): all 12 Table 1 metric values are defined on {[}0,
1{]}.

Weight normalization is critical: with 12 metrics and raw weights
averaging \textasciitilde0.6, the sum approaches 7.2. Without
normalization, weighted sums would saturate and collapse variance.

\section{Dynamic Thresholding and Homeostatic
Control}\label{sec-thresholding}

The write gate compares composite scores against an adaptive threshold
θ\_dynamic (written as theta\_dynamic when referring to the recorded
variable). This section details the threshold evolution algorithm
incorporating Bayesian prior-evidence blending and homeostatic rate
control.

\subsection{Prior-Evidence Blending}\label{prior-evidence-blending}

The threshold prior derives from knob settings:

\begin{verbatim}
θ_prior(F, S, T) = lerp(0.10, 0.30, T) × (1 − 0.3S)
\end{verbatim}

Observed evidence comes from the 90th percentile of recent scores:

\begin{verbatim}
recent_scores ← tail(score_stream, win_score(T))
if |recent_scores| == 0:
    observed_p90 ← θ_prior  # no evidence yet
else:
    observed_p90 ← percentile(recent_scores, 90)
\end{verbatim}

Prior and evidence masses weight the blend:

\begin{verbatim}
ρ_prior ← prior_mass(T) = round(lerp(2, 32, T))
ρ_obs ← u(t) × |recent_scores|
\end{verbatim}

The target threshold blends prior and evidence:

\begin{verbatim}
θ_target ← (ρ_prior × θ_prior + ρ_obs × observed_p90) /
            max(ε, ρ_prior + ρ_obs)
\end{verbatim}

High Stability increases prior mass, making the system more resistant to
observed deviations. High uncertainty increases evidence mass, allowing
faster adaptation to volatile conditions.

\subsection{Homeostatic Rate Control}\label{homeostatic-rate-control}

The controller maintains write rates near the target setpoint through
continuous-time estimation with effective sample size (ESS) reliability
weighting.

\subsubsection{Rate Estimation}\label{rate-estimation}

\begin{verbatim}
Δt ← now_s() − to_s(last_rate_timestamp)
dt_min(T) = 10^(−4 + 2T)  # stability-derived minimum Δt (seconds)
Δt ← max(Δt, dt_min(T))
τ_dt(T) = lerp(0.5, 2.0, T)
α_dt ← 1 − exp(−Δt / τ_dt(T))
dt_ema ← (1 − α_dt) × dt_ema + α_dt × Δt
dt_floor(T) = lerp(0.1, 0.5, T)  # stability-derived minimum cadence
dt_base ← max(dt_ema, dt_floor(T))
\end{verbatim}

The rate time constant scales with Stability:

\begin{verbatim}
τ_rate ← max(2^(3T) × dt_base, dt_floor(T))
α ← 1 − exp(−Δt / τ_rate)
\end{verbatim}

Instantaneous rate estimation with bias correction. Δwrites is the
binary indicator (0 or 1) of whether a write occurred during the current
timestep:

\begin{verbatim}
Δwrites ← 1 if write_memory else 0  # binary write event
\end{verbatim}

Normative note (MUST): Δwrites is computed from the current step's
write\_memory decision. Rate state updates occur after the write
decision and affect subsequent timesteps.

\begin{verbatim}
ρ_inst ← (Δwrites / Δt) × 60  # writes per minute
m_rate_t ← (1 − α) × m_rate + α × ρ_inst
denom ← max(1 − (1 − α)^(rate_ticks + 1), ε)
ρ_hat_t ← m_rate_t / denom  # bias-corrected estimate for next step
m_rate ← m_rate_t
rho_hat_prev ← ρ_hat_t
rate_ticks ← rate_ticks + 1
last_rate_timestamp ← now_ms()  # update after rate computation
\end{verbatim}

\subsubsection{Effective Sample Size}\label{effective-sample-size}

ESS estimates the effective number of independent samples in the EMA,
using a heuristic inspired by Liu and Chen (1998):

\begin{verbatim}
β ← max(0, 1 − α)
ESS_cap(T) = lerp(30, 120, T)
ESS ← min((1 + β) / max(1 − β, ε), ESS_cap(T))
reliability ← 1 − exp(−ESS × (1 − T))
\end{verbatim}

High Stability dampens reliability, preventing aggressive corrections in
conservative regimes.

\subsubsection{Homeostatic Correction}\label{homeostatic-correction}

The rate error drives threshold adjustment. The κ\_* gains are derived
from the knobs to avoid fixed behavioral constants:

Normative note (MUST): rho\_hat\_prev is stored state entering the
timestep. After the write decision, the rate-state update computes
ρ\_hat\_t; the assignment rho\_hat\_prev ← ρ\_hat\_t occurs at end of
step (see Appendix D).

\begin{verbatim}
rate_error ← tanh((rho_hat_prev − rate_target) /
                  max(rate_target, ε))
κ_r(F,S,T) = lerp(0.06, 0.14, S) × lerp(1.1, 0.9, T)
cap_homeo(F,S,T) = lerp(0.35, 0.15, T) × lerp(0.8, 1.2, S) × hysteresis
Δθ_homeo ← clamp(reliability × κ_r(F,S,T) × (1 − T) ×
                  (1 − maturity(t)) × rate_error,
                  −cap_homeo(F,S,T), +cap_homeo(F,S,T))
\end{verbatim}

The correction scales with reliability and is attenuated by both
Stability and maturity, ensuring conservative, mature systems make
minimal homeostatic adjustments.

\subsubsection{Sensitivity-Based Threshold
Adjustment}\label{sensitivity-based-threshold-adjustment}

Sensitivity modulates threshold based on recent score volatility:

\begin{verbatim}
recent_scores ← tail(score_stream, win_score(T))
if |recent_scores| < 2:
    σ_scores ← 0
else:
    σ_scores ← std(recent_scores)
κ_sens(F,S,T) = lerp(0.04, 0.12, S) × lerp(1.1, 0.9, T)
cap_sens(F,S,T) = lerp(0.30, 0.10, T) × lerp(0.9, 1.1, S) × hysteresis
σ_ref(S,T) = lerp(0.08, 0.14, S) × lerp(1.1, 0.9, T)
Δθ_sens ← clamp(−κ_sens(F,S,T) × S × (σ_scores − σ_ref(S,T)),
                −cap_sens(F,S,T), +cap_sens(F,S,T))
\end{verbatim}

High score variance with high Sensitivity lowers threshold, capturing
more volatile signals.

\subsubsection{Precision-Based Threshold
Adjustment}\label{precision-based-threshold-adjustment}

Focus-driven precision tightens threshold when structural coherence is
high:

\begin{verbatim}
κ_prec(F,S,T) = lerp(0.04, 0.10, F) × lerp(1.1, 0.9, T)
cap_prec(F,S,T) = lerp(0.25, 0.08, T) × lerp(0.9, 1.1, F) × hysteresis
Δθ_prec ← clamp(κ_prec(F,S,T) × F × (coherence_struct_t − 0.5),
                −cap_prec(F,S,T), +cap_prec(F,S,T))
\end{verbatim}

High structural coherence with high Focus raises threshold, enforcing
stricter relevance filtering.

\subsection{Threshold Integration}\label{threshold-integration}

All threshold deltas combine and pass through safety limiting. The
emotion and mood deltas from Section~\ref{sec-emotional-threshold} and
Section~\ref{sec-mood} are denoted Δθ\_emo = ΔThreshold\_emotion\_t and
Δθ\_mood = ΔThreshold\_mood\_t:

\begin{verbatim}
Δθ_total ← Δθ_sens + Δθ_homeo + Δθ_prec + Δθ_emo + Δθ_mood
cap_total ← max_ΔT_per_min(t) × (Δt / 60.0)
Δθ_limited ← clamp(Δθ_total, −cap_total, +cap_total)
θ_dynamic_t ← clamp(EWMA(θ_dynamic_{t−1}, θ_target,
                        α = α_T(t)) + Δθ_limited,
                    T_min(t), T_max(t))
theta_dynamic ← θ_dynamic_t
\end{verbatim}

Hysteresis evolves toward the stability-derived base:

\begin{verbatim}
hysteresis_t ← clamp(EWMA(hysteresis,
                     base_band(T), α = α_T(t)),
                band_min, band_max)
hysteresis ← hysteresis_t
\end{verbatim}

\subsection{Write Pacing and Memory
Accumulation}\label{sec-write-pacing}

The write gate operates per-signal, but coherent ``thoughts'' often span
multiple signals. This section introduces memory-level accumulation that
groups signals into natural units before storage decisions, inspired by
Event Segmentation Theory (Zacks \& Swallow, 2007).

This approach draws from EM-LLM (Fountas et al., 2024), which segments
token sequences into episodic events using surprise-based boundary
detection refined by graph-theoretic cohesion metrics. Their work shows
that combining prediction error signals with within-segment coherence
produces boundaries strongly correlated with human event perception. Our
adaptation uses EMA prediction error (surprisal\_t) for surprise and
cosine similarity for cohesion, while drift provides auxiliary boundary
pressure, enabling modality-agnostic operation across text tokens, audio
chunks, video frames, or any signal stream.

\subsubsection{Memory Accumulator State}\label{sec-accumulator}

Each source stream maintains accumulator state:

\begin{verbatim}
μ_acc ← 256d running mean embedding
drift_acc ← accumulated drift within memory
s_sum ← sum of signal scores in memory
s_max ← max signal score in memory
n ← count of signals in memory
e_peak ← embedding of highest-scoring signal
emo_max ← 0  # max emotion_intensity_t within the unit
arousal_sum ← 0  # sum of arousal_t within the unit (for avg)
t_start ← now_ms()  # timestamp of accumulation start (ms since epoch)
last_signal_ts ← now_ms()  # timestamp of previous signal (ms, for gap detection)
last_write_ts ← 0  # ms; 0 means "no prior write" for refractory
eta_acc ← 0  # drift EWMA state
coherence_prev ← 0  # previous coherence (initialize to 0)
acc_signals_window ← []  # ring buffer of recent embeddings for coherence
\end{verbatim}

Reset behavior: reset\_accumulator() clears μ\_acc, drift\_acc, s\_sum,
s\_max, n, e\_peak, emo\_max, arousal\_sum, eta\_acc, coherence\_prev,
and sets acc\_signals\_window ← {[}{]} (and refreshes
t\_start/last\_signal\_ts for the next unit), but retains
last\_write\_ts so the refractory term remains well-defined across
boundaries.

\begin{verbatim}
reset_accumulator(): acc_signals_window ← []  # MUST clear coherence window at boundaries
\end{verbatim}

On each signal, update running statistics (note: n is the count before
this signal):

\begin{verbatim}
signal_gap_s ← now_s() − to_s(last_signal_ts)  # compute BEFORE updating last_signal_ts

win_coh(T) = round(lerp(8, 32, T))  # coherence window size
acc_signals_window ← tail(acc_signals_window, win_coh(T))

μ_acc ← (n × μ_acc + x_t) / (n + 1)
n ← n + 1
drift_acc ← drift_acc + (drift_mag_t / 2)  # accumulate normalized context-centroid drift (drift_mag_t ∈ [0,2])
s_sum ← s_sum + score_t
emo_max ← max(emo_max, emotion_intensity_t)
arousal_sum ← arousal_sum + arousal_t
if score_t > s_max:
    s_max ← score_t; e_peak ← x_t
last_signal_ts ← now_ms()  # update timestamp for next gap calculation (ms)
\end{verbatim}

\subsubsection{Hybrid Drift and Coherence
Tracking}\label{sec-drift-coherence}

Boundary detection operates on the \textbf{current unflushed memory
centroid} rather than the raw signal, keeping segmentation tied to the
evolving thought unit:

\begin{verbatim}
ε_noise(T) = lerp(0.01, 0.05, 1 − T)  # stability-derived noise floor
# Use centroid drift (μ_acc) instead of per-signal drift
mu_prev ← mu_acc_{t−1}
mu_curr ← mu_acc_t
d_step ← max(1 − cos(mu_curr, mu_prev) − ε_noise(T), 0)
eta_prev ← eta_acc  # baseline before updating EWMA
\end{verbatim}

Memory coherence tracks similarity within the current memory
accumulation window (range {[}−1, 1{]} from mean cosine):

\begin{verbatim}
current_window ← acc_signals_window  # embeddings from the current unit before mu_curr
if |current_window| == 0:
    coherence_curr ← 1.0  # empty-window fallback
else:
    coherence_curr ← mean([cos(mu_curr, x_i) for x_i in current_window])
# After computing coherence_curr, append mu_curr for the next step
acc_signals_window.append(mu_curr); acc_signals_window ← tail(acc_signals_window, win_coh(T))
# coherence_prev is the stored coherence value from the previous step
\end{verbatim}

Note: coherence\_mem is distinct from coherence\_struct
(Section~\ref{sec-struct-coherence}). The former tracks within-memory
similarity using raw mean cosine, while the latter measures
variance-based integration with broader context. This dual-signal
approach mirrors EM-LLM's boundary detection mechanism. In their
formulation, boundaries occur where surprise (token-level prediction
error) exceeds a threshold and segment cohesion drops. In Cortext,
surprisal\_t provides the prediction-error signal, while drift spike
provides auxiliary boundary pressure and coherence\_mem drop captures
within-memory similarity degradation.

\subsubsection{Natural Boundary Detection}\label{sec-boundary}

Boundary score combines drift spike and coherence drop:

\begin{verbatim}
weight_drift_component = lerp(0.6, 0.4, T)  # weight on drift
weight_gap_component = lerp(0.05, 0.20, T)  # soft gap influence (low weight)
weight_drift_component = weight_drift_component × (1 − weight_gap_component)
weight_coh_component = 1 − weight_gap_component − weight_drift_component
ε0(T) = lerp(0.005, 0.02, 1 − T)  # cold-start guard for eta_prev
if eta_prev < ε0(T):
    drift_spike ← 0
else:
    drift_spike ← (d_step − eta_prev) / max(eta_prev, ε)
eta_acc ← EWMA(eta_prev, d_step, α = lerp(0.3, 0.1, T))
coh_drop01 ← clamp((coherence_prev − coherence_curr) / 2, 0, 1)
coherence_prev ← coherence_curr  # update for next step

# Adaptive gap signal (dynamic, not a hard trigger)
dt_ref ← max(dt_ema, dt_floor(T))  # EWMA of inter-arrival Δt in seconds (robust to jitter)
gap_scale(T) = lerp(3.0, 8.0, T)  # expected pause multiplier
gap_ref_s ← gap_scale(T) × dt_ref
gap_z ← (signal_gap_s − gap_ref_s) / max(gap_ref_s, ε)
gap_score ← sigmoid(gap_z)

boundary_score ← weight_drift_component × sigmoid(drift_spike) +
                  weight_coh_component × coh_drop01 +
                  weight_gap_component × gap_score
boundary_score ← clamp(boundary_score, 0, 1)
\end{verbatim}

Boundary threshold and limits:

\begin{verbatim}
b_thresh(F, S) = lerp(0.4, 0.7, F) × lerp(1.1, 0.9, S)
max_mem_time(T) = lerp(30, 120, T)  # seconds
max_mem_drift(S) = lerp(0.8, 2.0, S)  # cumulative drift cap
\end{verbatim}

Trigger memory flush when:

\begin{verbatim}
mem_elapsed ← now_s() − to_s(t_start)
should_flush = (boundary_score > b_thresh(F, S)) OR
               (mem_elapsed > max_mem_time(T)) OR
               (drift_acc > max_mem_drift(S))
\end{verbatim}

where signal\_gap\_s = now\_s() − to\_s(last\_signal\_ts) is computed at
the start of signal processing (before last\_signal\_ts is updated). Gap
timing influences boundary\_score via gap\_score rather than forcing an
unconditional flush, which avoids splitting mid-thought when upstream
processing introduces non-deterministic latency.

\subsubsection{Spike Bypass (Flashbulb
Flush)}\label{spike-bypass-flashbulb-flush}

High-salience signals bypass accumulation and flush immediately,
capturing preceding context as a coherent memory unit:

\begin{verbatim}
spike_margin(S,T) = lerp(0.2, 0.5, T) × lerp(1.2, 0.8, S)  # above θ_dynamic
spike_bypass = score_t > (θ_dynamic + spike_margin(S,T))
force_write ← false
\end{verbatim}

When spike\_bypass triggers:

\begin{verbatim}
if spike_bypass:
    should_flush = true   # force a boundary
    force_write = true   # bypass S_window > θ_memory
\end{verbatim}

This ensures flashbulb moments capture their surrounding context rather
than creating isolated micro-memories.

\subsubsection{Window Score and
Refractory}\label{window-score-and-refractory}

Memory-level score combines peak and average with coverage bonus:

\begin{verbatim}
s_avg ← s_sum / max(n, 1)
α(F) = lerp(0.3, 0.7, F)  # peak vs avg weight
coverage ← min(n / n_ctx(T), 1.0)  # memory completeness
β(S) = lerp(0.05, 0.15, S)  # coverage weight
S_window ← α(F) × s_max + (1 − α(F)) × s_avg + β(S) × coverage
\end{verbatim}

Write refractory suppresses rapid successive writes:

\begin{verbatim}
τ_write_refrac(T) = lerp(5, 30, T)  # seconds
k_write_refrac = lerp(0.3, 0.1, T)
Δt_write ← now_s() − to_s(last_write_ts)
M_write_refrac ← 1.0 + k_write_refrac × exp(−Δt_write / τ_write_refrac(T))
\end{verbatim}

Final write decision:

\begin{verbatim}
θ_memory ← θ_dynamic × M_write_refrac
write_memory = force_write OR (should_flush AND (S_window > θ_memory))
if write_memory: last_write_ts ← now_ms()  # update refractory timestamp (ms)
\end{verbatim}

Normative rule (MUST): if should\_flush is true, the current unit must
be finalized. If write\_memory is false, discard the unit and
reset\_accumulator() anyway (do not update last\_write\_ts). This
prevents perpetual should\_flush states (time cap / drift cap / gap cap)
while never resetting.

Trace note: if a run trace reports both write\_decision and stored,
interpret write\_decision as the boolean gate outcome at the boundary
(the write\_memory predicate above). Interpret stored as the eventual
recording outcome (after any final safety checks).

Representative embedding blends accumulator mean with peak:

\subsubsection{Representative Embedding}\label{sec-rep-embedding}

\begin{verbatim}
ρ(F) = lerp(0.3, 0.7, F)  # mean vs peak blend
e_rep ← l2_normalize(ρ(F) × μ_acc + (1 − ρ(F)) × e_peak)
\end{verbatim}

On write: store e\_rep with metadata \{n, s\_max, s\_avg, drift\_acc,
mem\_elapsed, s\_emotion\_max=emo\_max, s\_arousal\_avg=arousal\_sum /
max(n, 1)\}. Append e\_rep to memory\_stream and
recent\_memory\_centroids. Reset accumulator for next unit.

\section{Reinforcement and Decay Dynamics}\label{sec-reinforcement}

\subsection{Memory Strength Model}\label{sec-memory-strength}

Each memory m maintains a strength value updated through use-frequency
tracking and exponential decay. The update uses a sensitivity-modulated
learning rate and a binary usage indicator:

\begin{verbatim}
α_min_S = 0.05; α_span_S = 0.35
α_S(t) = α_min_S + S × α_span_S × u(t)
used_flag(m) = 1 if m was retrieved and used in current step, else 0

use_frequency_t ← EWMA(use_frequency_{t−1},
                       used_flag(m), α = α_S(t))
λ_t ← ln(2) / half_life_t
strength_t ← clamp(strength_{t−1} × exp(−λ_t × Δt) +
                   S × use_frequency_t, 0, 1)
\end{verbatim}

Δt is measured per memory using its last access timestamp:

\begin{verbatim}
Δt(m) ← now_s() − to_s(m.last_access); if last_access is unset, use created_at
\end{verbatim}

Memories falling below the periphery cutoff are candidates for eviction:

\begin{verbatim}
periphery_cutoff(T) = lerp(0.05, 0.25, T)
if strength_t < periphery_cutoff(T): evict(m)
\end{verbatim}

\subsection{Contextual Gain and Per-Memory Stability
(Definitions)}\label{contextual-gain-and-per-memory-stability-definitions}

Contextual gain is a per-retrieval signal used by multiple feedback
loops:

\begin{verbatim}
contextual_gain_t(m) ← cos(x_t, m.embedding)    # if m was retrieved and used (range [−1, 1])
contextual_gain_t(m) ← undefined                # otherwise
\end{verbatim}

When undefined, treat contextual\_gain\_t(m) as 0 in downstream updates.
Persisted per-memory gain is tracked as a smoothed value:

\begin{verbatim}
L_cg = round(lerp(8, 32, T))
α_cg = 2 / (L_cg + 1)
contextual_gain(m) ← EWMA(contextual_gain(m), contextual_gain_t(m), α_cg)
\end{verbatim}

Per-memory stability initializes at creation and is bounded:

\begin{verbatim}
stability(m)_init = 1.0
stability(m) ← clamp(stability(m), 0.0, 2.0)
\end{verbatim}

Retrieved vs used:

\begin{verbatim}
retrieved(m) = memory returned by retrieval (vector or graph expansion)
used(m) = retrieved(m) AND injected into active context after gate decisions
\end{verbatim}

Contextual gain is only measured for used memories.

\subsection{Influence-Weighted
Updates}\label{influence-weighted-updates}

When contextual gain signals are available, influence factors modulate
reinforcement:

\begin{verbatim}
influence_factor ← (used_count / max(retrieved_count, 1)) ×
                    clamp(contextual_gain(m), −1, +1)
strength_t ← clamp(strength_{t−1} × exp(−λ_t × Δt) +
                   S × use_frequency_t + F × influence_factor, 0, 1)
\end{verbatim}

\subsection{Causal Feedback Loop}\label{sec-feedback}

The system tracks causal influence of retrieved memories on generation
quality through contextual gain---the semantic alignment (cosine
similarity) between the current signal embedding and the retrieved
memory embedding.

\subsubsection{Focus Feedback}\label{focus-feedback}

\begin{verbatim}
αF_base = 0.10; βF_base = 0.05
weight_relevance_t ← weight_relevance
attention_width_t ← attention_width
for each used memory m:
    if contextual_gain(m) > 0:
        weight_relevance_t += αF_base × contextual_gain(m)
        attention_width_t *= (1 − βF_base)
    else:
        attention_width_t *= (1 + βF_base)
weight_relevance ← clamp(weight_relevance_t, 0, 1)
attention_width_t ← clamp(attention_width_t,
                          attention_width_min, attention_width_max)
attention_width ← attention_width_t
\end{verbatim}

Positive contextual gain narrows attention and boosts relevance
weighting; negative gain widens attention to explore alternatives.

\subsubsection{Sensitivity Feedback}\label{sensitivity-feedback}

\begin{verbatim}
η_base = 0.10
weight_novelty_t ← weight_novelty
for each used memory m:
    novelty_reward ← 1 − sim(m.embedding, recent_context)
    weight_novelty_t += η_base × (novelty_reward ×
                         contextual_gain(m) −
                         redundancy(m, recent_context))
weight_novelty_t ← clamp(weight_novelty_t, 0, 1)
weight_novelty ← weight_novelty_t
\end{verbatim}

This rewards novelty that proves useful while penalizing redundant
retrievals.

\subsubsection{Stability Feedback}\label{stability-feedback}

\begin{verbatim}
γT_base = 0.05
for each used memory m:
    if contextual_gain(m) > 0:
        stability_t(m) ← stability(m) + γT_base
    else:
        stability_t(m) ← stability(m) × (1 − γT_base)
    stability(m) ← clamp(stability_t(m), 0.0, 2.0)
\end{verbatim}

The mean stability of used memories provides adjustment to the half-life
target:

\begin{verbatim}
adj ← clamp(mean(stability(m_used)) − 1.0, −0.25, +0.25)
ΔHalfLife_adj_t ← adj
\end{verbatim}

This factor is consumed by the Stability update
(Section~\ref{sec-stability-update}), avoiding conflicting adjustments
between feedback mechanisms.

\subsection{Generation Influence
Tracking}\label{generation-influence-tracking}

When generation embeddings are available, influence incorporates output
trajectory:

\begin{verbatim}
Δḡ ← l2_normalize(ḡ_t) − l2_normalize(ḡ_{t−1})
drift_mag_gen ← ‖Δḡ‖
drift_contribution(m) ← (drift_mag_gen / 2) ×
                         max(0, cos(m.embedding, l2_normalize(Δḡ)))
\end{verbatim}

Total influence blends contextual gain, generation similarity, and drift
contribution:

\begin{verbatim}
λ₁ = 0.5; λ₂ = 0.4; λ₃ = 0.3
influence(m) ← λ₁ × contextual_gain(m) +
                λ₂ × cos(m.embedding, ḡ_t) −
                λ₃ × drift_contribution(m)
\end{verbatim}

Sustained influence accumulates over a stability-dependent horizon:

\begin{verbatim}
L_sustain(T) = round(lerp(3, 5, T))
sustained_influence ← EWMA(sustained_influence,
                           influence(m),
                           α = 2 / (L_sustain(T) + 1))
\end{verbatim}

\section{Advanced Cognitive Processes}\label{sec-advanced}

This section presents algorithms modeling higher-order cognitive
phenomena: working memory maintenance, metacognitive monitoring,
reconsolidation dynamics, and serial position effects.

\subsection{Working Memory Gates}\label{working-memory-gates}

Following Cowan (2001) capacity constraints, working memory maintains a
limited number of active items. Working memory holds coherent memories
as defined in Section~\ref{sec-write-pacing}, preserving the full
content and signal sequence:

\begin{verbatim}
base_capacity = round(lerp(5, 3, S) + lerp(−1, 1, F))
\end{verbatim}

This yields a range of approximately 2-6 memories, broadening the 4±1
chunk limit to accommodate task-dependent requirements. High Sensitivity
reduces capacity (faster turnover), while high Focus modulates breadth.

\subsubsection{Active Memory Structure}\label{active-memory-structure}

Each active memory holds a coherent memory with its full content and
metadata:

\begin{verbatim}
memory.content ← concatenated signal blobs
memory.source_id ← source identifier (e.g., 'user', 'assistant')
memory.modality ← primary modality ('text', 'audio', 'image')
memory.blob_ids ← [blob_1, blob_2, ..., blob_n]  # blob refs
memory.embedding ← e_rep  # representative embedding (@sec-rep-embedding)
memory.signals ← [x_1, x_2, ..., x_n]  # ordered signal embeddings
memory.metadata ← {n, s_max, s_avg, drift_acc, mem_elapsed,
                   s_emotion_max, s_arousal_avg}
\end{verbatim}

The emotional metrics (s\_emotion\_max, s\_arousal\_avg) are accumulated
during memory formation for use by Emotional Consolidation
(Section~\ref{sec-emotional-consolidation}).

\subsubsection{Maintenance Cost}\label{maintenance-cost}

Maintenance incurs cognitive cost:

\begin{verbatim}
maintenance_cost_per_memory = lerp(0.05, 0.15, S)
complexity_penalty = manifold_complexity × lerp(0.5, 1.5, S)
\end{verbatim}

The manifold\_complexity is defined as a normalized local variability
proxy (see Appendix B): manifold\_complexity ← clamp((1 −
mean\_cos\_window) / 2, 0, 1).

\subsubsection{Memory-Level Gating}\label{memory-level-gating}

Gating thresholds determine entry and chunking:

\begin{verbatim}
chunking_threshold = lerp(0.7, 0.9, F)
gate_threshold = lerp(0.1, 0.4, F)
rehearsal_rate = lerp(0.5, 2.0, S)
memory_dedication_strength = lerp(0.3, 0.9, T)
\end{verbatim}

Working memory gating evaluates coherent memories at accumulation
boundaries (Section~\ref{sec-boundary}), not individual signals:

\begin{verbatim}
on_memory_boundary:
    [α, β, γ] ← normalize([lerp(0.55, 0.70, F),   # window score weight
                        lerp(0.20, 0.35, F),   # task relevance weight
                        lerp(0.10, 0.30, S)])   # novelty-to-WM weight
    memory_benefit ← α × S_window + β × relevance_to_task(μ_acc, task_context) +
                       γ × novelty_to_set(μ_acc, {m.embedding | m ∈ active_memories})
    margin ← memory_benefit − gate_threshold
    k ← |active_memories|
    C ← max(base_capacity, 1)
    p_cap ← 3
    capacity_pressure(k, C) ← 1 + max(0, (k − C) / C)^p_cap
    raw_cost ← (maintenance_cost_per_memory × k + complexity_penalty) ×
               capacity_pressure(k, C)
    total_cost ← raw_cost / (1 + raw_cost)  # squash to [0, 1)
    accept_memory = (margin ≥ total_cost)
\end{verbatim}

Normative note (MUST): total\_cost is normalized to {[}0, 1{]} to keep
gating on the same score scale as margin. This avoids cold-start or
capacity regimes where costs can exceed 1 and make acceptance
vanishingly rare.

Note that total\_cost is computed from existing active memories only (k
is the current count), so k=0 yields no bootstrap penalty. Capacity
pressure activates only when k \textgreater{} C.

\subsubsection{Chunking at Memory Level}\label{chunking-at-memory-level}

Chunking operates on memory embeddings to merge related content from the
same source:

\begin{verbatim}
similar_memories ← {m ∈ active_memories | cos(m.embedding, memory.e_rep) > chunking_threshold AND
                              m.source_id == memory.source_id}
if |similar_memories| > 0:
    merge_into_chunk(similar_memories, memory)
\end{verbatim}

\subsection{Metacognitive Monitoring}\label{metacognitive-monitoring}

The system implements feeling-of-knowing (FOK) and tip-of-tongue (TOT)
detection following Hart (1965) framework:

\begin{verbatim}
FOK_threshold = lerp(0.2, 0.5, F)
TOT_detection = (FOK > lerp(0.5, 0.8, F)) AND
                 (retrieval_strength < lerp(0.4, 0.2, F))
\end{verbatim}

TOT occurs when metacognitive confidence is high but retrieval strength
is low---the characteristic experience of knowing one knows something
but being unable to access it.

Additional metacognitive parameters:

\begin{verbatim}
confidence_decay_rate = lerp(0.01, 0.1, 1 − T)
unknown_threshold = lerp(0.3, 0.1, F)
strategy_switch_latency = lerp(500, 100, S)  # ms
certainty_requirement = lerp(0.6, 0.9, T)
metacognitive_sensitivity = F × (1 + 0.5 × S)
\end{verbatim}

\subsection{Memory Reconsolidation}\label{memory-reconsolidation}

Following Nader, Schafe, and Le Doux (2000), retrieved memories enter a
labile state permitting modification:

\begin{verbatim}
τ_labile = lerp(30, 300, T)  # seconds
reconsolidation_gain = lerp(0.2, 0.02, T)
lability_susceptibility = (1 − T) × (0.5 + 0.5 × S)
\end{verbatim}

During the lability window, memories can drift toward current context:

\begin{verbatim}
drift_magnitude = (1 − T) × S × lability ×
                   contextual_relevance
\end{verbatim}

Reconsolidation effects propagate to semantically related memories with
decay:

\begin{verbatim}
ripple_decay = lerp(0.5, 0.1, T)  # per semantic hop
\end{verbatim}

\subsection{Retrieval Competition}\label{retrieval-competition}

Retrieved memories compete through lateral inhibition, modeling
retrieval-induced forgetting (Anderson, Bjork, and Bjork 1994):

\begin{verbatim}
inhibition_radius = lerp(0.5, 0.85, F)
winners_k = round(lerp(7, 3, F))
suppression_per_retrieval = lerp(0.1, 0.01, T) ×
                             (1 − winning_activation)
recovery_time_RIF = lerp(300, 1800, T)  # seconds
\end{verbatim}

High Focus produces narrow winner-take-all dynamics; low Focus permits
broader activation.

\subsection{Predictive Pre-activation}\label{predictive-pre-activation}

The system pre-activates memories predicted to be relevant based on the
EMA expectation state (\texttt{x\_pred\_ema}) and recent context:

\begin{verbatim}
prediction_horizon = round(lerp(2, 8, F))
pre_activation_decay = lerp(0.7, 0.3, T)
prediction_conf_threshold = lerp(0.3, 0.7, F)
surprise_sensitivity = S × lerp(2.0, 0.5, T)
\end{verbatim}

When predictions fail (high surprise), the system increases the refresh
rate of pre-activation:

\begin{verbatim}
update_rate_on_surprise = lerp(0.2, 0.02, T) × S
\end{verbatim}

\subsection{Serial Position Effects}\label{serial-position-effects}

The architecture models primacy, recency, and distinctiveness effects
observed in human memory (Murdock Jr 1962):

\begin{verbatim}
primacy_window = round(lerp(5, 2, F))
primacy_bonus = lerp(1.2, 2.0, S)
recency_window = round(lerp(7, 3, F))
rehearsal_curve_depth = lerp(0.2, 0.6, S)
\end{verbatim}

The von Restorff (isolation) effect enhances memory for distinctive
items (Hunt 1995):

\begin{verbatim}
distinctiveness_threshold = lerp(0.6, 0.8, F)
von_restorff_multiplier = lerp(1.5, 3.0, S)
\end{verbatim}

Items in the middle region suffer interference:

\begin{verbatim}
interference_zone = positions[primacy_window+1 : −recency_window]
middle_suppression = lerp(0.8, 0.5, S) × (1 − F)
\end{verbatim}

\subsection{Emotional Consolidation}\label{sec-emotional-consolidation}

High-emotion events trigger enhanced consolidation, following McGaugh
(2004) findings. As detailed in Section~\ref{sec-activity},
consolidation operates on stored memory metadata:

\begin{verbatim}
θ_intensity = lerp(0.6, 0.8, 1 − S)
θ_arousal = lerp(0.4, 0.2, S)
# Consolidation uses stored memory emotional metadata
trigger = (m.metadata.s_emotion_max ≥ θ_intensity) AND
           (m.metadata.s_arousal_avg ≥ θ_arousal)
\end{verbatim}

Flashbulb memories receive extended half-life bonuses based on the
memory's peak emotional intensity:

\begin{verbatim}
flashbulb_threshold = lerp(0.9, 0.4, S)
# Half-life bonus uses stored memory emotional peak
emotional_half_life_bonus = exp(lerp(0, ln(3), S)) ×
                             (1 + m.metadata.s_emotion_max)
\end{verbatim}

The emotional metrics (s\_emotion\_max, s\_arousal\_avg) are accumulated
during memory formation and stored with the memory
(Section~\ref{sec-accumulator}).

\begin{verbatim}
cascade_radius = round(lerp(1, 5, S))
cascade_decay = lerp(0.7, 0.3, S)
\end{verbatim}

\section{Consolidation and Graph Integration}\label{sec-consolidation}

The consolidation system transforms episodic memories into semantic
structures through clustering, summarization, and knowledge graph
construction. Memory-level storage (Section~\ref{sec-write-pacing})
ensures each embedding represents a coherent unit.

\subsection{Consolidation Triggers}\label{sec-consolidation-triggers}

Consolidation operates on stored memory representatives (e\_rep from
Section~\ref{sec-rep-embedding}), not individual signals. It activates
under capacity, rate, or temporal conditions:

\begin{verbatim}
should_consolidate = (memory_count > consolidation_threshold) OR
                      (m_rate < rate_target / 2) OR
                      (elapsed_time > consolidation_interval)
\end{verbatim}

Definitions (knob-derived):

\begin{verbatim}
consolidation_interval(T) = lerp(300, 3600, T)  # seconds (5 min → 1 hour)
consolidation_threshold(T) = n_ctx(T) × win_score(T)  # memory-count trigger
merge_threshold(F) = lerp(0.85, 0.95, F)  # similarity required for clustering
\end{verbatim}

The consolidation rate adapts to Stability and Sensitivity (write\_rate
tracks memory writes, not signal writes):

\begin{verbatim}
rate_consolidate = (1 / max(consolidation_interval, 1)) ×
                    (0.3 + 0.7T) × (1 − 0.5S)
\end{verbatim}

\subsubsection{Activity-Aware Scheduling}\label{sec-activity}

Consolidation runs during idle periods and preempts for retrieval. The
is\_accumulating\_memory check ensures consolidation doesn't interrupt
mid-memory accumulation:

\begin{verbatim}
idle_required(T) = round(0.25 × win_rate_s(T))
idle_for_s = now_s() − to_s(last_retrieval_ts)
# Consolidation waits for memory completion, not just signal arrival
should_start = (NOT is_accumulating_memory) AND
                (retrieval_queue_depth == 0) AND
                (idle_for_s ≥ idle_required(T))
\end{verbatim}

Consolidation begins only after the current memory has been flushed
(Section~\ref{sec-boundary}) and the idle period has elapsed. On
retrieval events, consolidation pauses, commits micro-batches, and
resumes when idle.

\subsection{Consolidation Scoring}\label{consolidation-scoring}

Each stored embedding represents a memory representative (e\_rep from
Section~\ref{sec-rep-embedding}). Each memory receives a consolidation
score determining merge priority:

\begin{verbatim}
score_consolidate(m) = weight_strength × strength(m) −
                        weight_redundancy × redundancy(m) +
                        weight_connectivity × connectivity(m) +
                        weight_stability × stability(m)
\end{verbatim}

Weights derive from knobs:

\begin{verbatim}
weight_strength = T; weight_redundancy = F; weight_connectivity = S; weight_stability = T
\end{verbatim}

Low-scoring memories are marked for merging.

\subsection{Clustering and
Summarization}\label{clustering-and-summarization}

Marked memories cluster via density-based methods (e.g., DBSCAN) or
k-means using embedding similarity:

\begin{verbatim}
cluster_i = {m_j | cos(m_j, μ_i) > merge_threshold}
μ_i = centroid(cluster_i)
\end{verbatim}

Summary nodes replace clusters:

\begin{verbatim}
summary.embedding = μ_i
summary.blob = summarize(fetch_blobs(cluster_i))
summary.metadata.sources = [m.id for m in cluster_i]
\end{verbatim}

\subsection{Semantic Extraction}\label{semantic-extraction}

For sufficiently large clusters, semantic extraction identifies labels
and relations:

\begin{verbatim}
extraction_batch_size = round(lerp(8, 32, T))
min_cluster_size = round(lerp(3, 10, F))
label_frequency_threshold = round(lerp(5, 15, T))
extraction_interval = lerp(300, 3600, T)  # 5 min → 1 hour
max_extractions_per_cycle = round(lerp(20, 5, T))
\end{verbatim}

Extraction uses structured prompting to identify labels/tags (people,
places, organizations, concepts) and relationships (co-occurrence,
implication, contradiction). Labels are stored as surface strings only
(no type/category field).

Label salience is derived from embeddings rather than model guesses:

\begin{verbatim}
e_label = encode(label_text)
salience(label) = clamp((cos(e_label, summary.embedding) + 1) / 2, 0, 1)
\end{verbatim}

If a label already exists, its stored salience is updated with
\texttt{s\_max\ =\ max(s\_max,\ salience(label))}.

\subsection{Knowledge Graph
Construction}\label{knowledge-graph-construction}

The graph comprises two node types:

\begin{itemize}
\tightlist
\item
  \textbf{Memory Nodes:} Summarized embeddings from merged clusters
\item
  \textbf{Label Nodes:} Extracted labels/tags derived from content blobs
\end{itemize}

Edge types capture relationships:

Edge weight convention (MUST): all association weights are normalized to
{[}0, 1{]}. For cosine-based edges, store weight01 = clamp((cos\_sim +
1) / 2, 0, 1).

\begin{itemize}
\tightlist
\item
  \textbf{co\_occurs:} Shared context or temporal proximity
\item
  \textbf{implies:} Directional correlation in embedding drift
\item
  \textbf{causes:} Directional correlation in embedding drift
\item
  \textbf{contradicts:} Strong negative similarity or constraint
  violation
\item
  \textbf{reinforces:} Frequent joint retrieval
\item
  \textbf{derived\_from:} Links summaries to source memories
\item
  \textbf{similar\_to:} High cosine similarity (soft equivalence)
\item
  \textbf{has\_label:} Attaches an extracted label/tag to a node
\end{itemize}

\subsubsection{Edge Construction}\label{edge-construction}

Co-occurrence edges derive from embedding similarity:

\begin{verbatim}
for (m_i, m_j) in cluster.sources:
    cos_sim ← cos(m_i.embedding, m_j.embedding)
    if cos_sim > lerp(0.85, 0.95, F):
        weight01 ← clamp((cos_sim + 1) / 2, 0, 1)
        create_edge(m_i, m_j, 'co_occurs', weight01)
\end{verbatim}

Causal edges derive from temporal drift:

\begin{verbatim}
temporal_order ← sort_by_timestamp(cluster.sources)
for i in range(len(temporal_order) − 1):
    m_i, m_j ← temporal_order[i], temporal_order[i+1]
    drift_vec ← m_j.embedding − m_i.embedding
    drift_mag ← ‖drift_vec‖
    if drift_mag > lerp(0.15, 0.35, T):
        weight01 ← clamp(drift_mag / 2, 0, 1)  # drift_mag ∈ [0,2] when embeddings are L2-normalized
        create_edge(m_i, m_j, 'causes', weight01)
\end{verbatim}

\subsection{Graph-Augmented Retrieval}\label{graph-augmented-retrieval}

Retrieval combines vector similarity with graph expansion. The query
vector is the \textbf{current accumulator centroid} (μ\_acc), which
aggregates the embeddings of the signals observed so far in the
unflushed memory. This makes retrieval responsive to the evolving
thought unit rather than any single micro‑signal.

\begin{verbatim}
# Query and re-rank using current memory centroid
q ← μ_acc  # memory centroid from Section 4.4.1
if |memory_stream| == 0:
    return []  # cold-start fallback (no retrieval candidates)
results_vec ← topK(vector_search(q, k=kNN_size(F)))
seed_nodes ← [r.id for r in results_vec]
if |seed_nodes| == 0 OR graph is empty:
    return results_vec  # deterministic fallback, skip expansion
expanded_nodes ← graph.traverse(seed_nodes, depth=graph_depth(T), min_edge_weight=min_edge_weight(F))
combined ← union(seed_nodes, expanded_nodes)
re_ranked ← sort_by(cos(q, embeddings(combined)))  # re-rank by memory centroid
# Exclude current-unit writes and WM overlaps before returning
dup_thresh = lerp(0.96, 0.88, F) × (0.98 + 0.02T)
eligible ← {m ∈ re_ranked | m.created_at < write_exclusion_ts}
eligible ← {m ∈ eligible | max_{w∈WM} cos(m.embedding, w) < dup_thresh}
return eligible
\end{verbatim}

Graph expansion uses recursive traversal with depth limits to find
related context that pure vector search might miss. Using the memory
centroid maintains consistency between vector search and graph expansion
results.

\section{Interrupt Gate and Streaming Integration}\label{sec-interrupt}

The interrupt gate controls when retrieved memories enter active context
during streaming generation. The gate balances novelty value against
disruption cost. The interrupt gate operates on memory-level context,
using centroids (μ\_acc) rather than individual signal embeddings for
novelty and relevance computation. All thresholds and gains in this
section are derived from the three knobs (F, S, T); no fixed behavioral
constants are introduced.

\subsection{Marginal Utility
Computation}\label{marginal-utility-computation}

Novelty thresholds scale with knobs and refractory state:

\begin{verbatim}
τ_novelty = lerp(0.10, 0.35, F) × (1 − 0.15S) × (1 + 0.3T)
τ_mu = lerp(0.08, 0.18, F) × (1 − 0.4S) × (1 + 0.4T)
retrieval_thresh(F) = lerp(0.25, 0.60, F)
\end{verbatim}

Refractory dynamics suppress rapid successive interrupts:

Interrupt state (per stream):

\begin{verbatim}
prev_x is unset on the first signal of a stream; set prev_x ← x_t after processing each signal
first_step ← (prev_x is unset for this stream)
if first_step:
    # cold start: do not update drift_accum
    drift_accum ← drift_accum
else:
    drift_accum ← drift_accum + cosine_dist(x_t, prev_x)  # cumulative drift
Δ ← drift_accum − drift_at_last_interrupt  # cumulative drift since last interrupt
τ_refrac = lerp(24, 96, T) × lerp(1.4, 1.0, S)
k_refrac = lerp(0.20, 0.05, T) × lerp(0.8, 1.2, F)
M_refrac = 1.0 + k_refrac × exp(−Δ / τ_refrac)
\end{verbatim}

On interrupt: set drift\_at\_last\_interrupt ← drift\_accum (resetting Δ
to 0 for subsequent signals).

Effective thresholds incorporate refractory pressure:

\begin{verbatim}
τ_novelty_eff = τ_novelty × M_refrac
τ_mu_eff = τ_mu × M_refrac
\end{verbatim}

\subsection{Marginal Utility Score}\label{marginal-utility-score}

The marginal utility (MU) of a candidate memory combines five factors.
Context comparisons use memory centroids rather than individual signal
embeddings:

\begin{verbatim}
# Context window contains recent memory centroids, not individual signals
ctx_window ← recent_memory_centroids  # bounded deque of recent memory representatives (e_rep)
q_retrieval ← μ_acc  # accumulator centroid query
included_set ← {embedding(m) | m already injected into the current context window}
\end{verbatim}

Fallback: if included\_set is empty, treat redundancy(·, included\_set)
= 0 and overlap\_star = −1.

\begin{verbatim}
if |ctx_window| == 0:
    ctx_centroid ← q_retrieval  # fallback: use current unit centroid
else:
    ctx_centroid ← mean(ctx_window)  # centroid of recent memory centroids

weights_mu_raw = [lerp(0.40, 0.60, F),   # coverage gain
         lerp(0.35, 0.25, F),   # relevance
         lerp(0.15, 0.25, S),   # redundancy penalty
         lerp(0.15, 0.25, S),   # incoherence penalty
         lerp(0.20, 0.50, S)]   # surprise bonus
[weight_cov, weight_rel, weight_red, weight_incoh, weight_surp] = normalize(weights_mu_raw)

mu = weight_cov × coverage_gain(candidate | included_set) +
      weight_rel × map01(cos(candidate, ctx_centroid)) −
      weight_red × redundancy(candidate, included_set) −
      weight_incoh × (1 − coherence_struct_t) +
      weight_surp × surprisal_t  # surprise bonus
\end{verbatim}

With map01 applied, each MU term is in {[}0, 1{]}, so μ is calibrated to
the same range as τ\_mu.

\subsection{Gate Decision Logic}\label{gate-decision-logic}

Duplicate suppression threshold:

\begin{verbatim}
dup_thresh = lerp(0.96, 0.88, F) × (0.98 + 0.02T)
K = round(lerp(10, 6, F))  # candidates to evaluate
\end{verbatim}

Higher Focus lowers \texttt{dup\_thresh}, making duplicate suppression
stricter.

\subsubsection{Write Exclusion Filter}\label{write-exclusion-filter}

Memories stored during the current accumulation unit are excluded from
interrupt consideration to prevent self-triggering. Using the
accumulation start timestamp ensures all memories written within the
current unit are excluded:

\begin{verbatim}
# Exclude memories written during current accumulation to prevent self-triggering
    write_exclusion_ts ← t_start  # start timestamp from Accumulator State section
candidates_eligible ← {c ∈ candidates | c.created_at < write_exclusion_ts}
\end{verbatim}

This filter is applied before novelty and marginal utility evaluation.
All subsequent gate logic operates on candidates\_eligible rather than
the raw candidate set, preventing recursive triggering within a coherent
thought unit.

\textbf{Working memory exclusion (normative):} retrieval must also
exclude candidates that overlap with active working‑memory slots. This
keeps ``what I just said'' in working memory rather than re‑injecting it
from long‑term memory. The exclusion uses the same duplication threshold
as the gate:

\begin{verbatim}
dup_thresh = lerp(0.96, 0.88, F) × (0.98 + 0.02T)
candidates_eligible ← {c ∈ candidates_eligible | max_{w∈WM} cos(c, w) < dup_thresh}
\end{verbatim}

Boundary-aware override permits lower-threshold interrupts at natural
boundaries:

\begin{verbatim}
boundary_mult = lerp(1.3, 2.0, F) × lerp(1.1, 0.9, S)
\end{verbatim}

The gate permits interrupt when:

\begin{verbatim}
at_drift_boundary = should_flush  # boundary signal from the current accumulator
if |candidates_eligible| == 0:
    allow_interrupt = false  # cold-start / empty-store fallback
else:
    candidate_star = argmax_{c ∈ candidates_eligible} mu(c)
    mu_star = mu(candidate_star)
    rel_star = map01(cos(candidate_star, ctx_centroid))
    if |included_set| == 0:
        overlap_star = −1.0
    else:
        overlap_star = max_{y ∈ included_set} cos(candidate_star, y)
    if |ctx_window| == 0:
        novelty_star = 1.0
    else:
        max_cos = max_{c ∈ ctx_window} cos(candidate_star, c)  # in [−1, 1]
        novelty_star = clamp((1 − max_cos) / 2, 0, 1)
    allow_interrupt =
        (rel_star ≥ retrieval_thresh(F)) AND
        (novelty_star ≥ τ_novelty_eff OR mu_star ≥ τ_mu_eff) AND
        (overlap_star < dup_thresh) AND
        (at_drift_boundary OR mu_star ≥ boundary_mult × τ_mu_eff)
\end{verbatim}

This logic suppresses low-drift interrupts unless the marginal utility
substantially exceeds threshold, while permitting normal-threshold
interrupts at natural transition points.

\subsection{Streaming Pacing}\label{streaming-pacing}

Streaming retrieval is gated by cumulative drift rate within the
accumulation unit. Retrieval checks trigger when drift exceeds threshold
or at boundaries. Retrieval uses q\_retrieval (the accumulator centroid)
captured before any accumulator reset for this step. Retrieval returns a
candidate pool already filtered by write‑exclusion and WM‑overlap rules;
the interrupt gate then applies novelty/utility thresholds and
redundancy penalties.

\begin{verbatim}
# Pacing tracks drift within current memory formation
\end{verbatim}

where cosine\_dist(u, v) = 1 − cos(u, v).

\begin{verbatim}
first_step ← (x_last_check is unset for this stream)  # MUST occur once per stream
if first_step: x_last_check ← x_t; drift_acc_pacing ← 0
drift_acc_pacing += cosine_dist(x_t, x_last_check)
pacing_thresh(S) = lerp(0.5, 0.1, S)
max_wait_drift(F) = lerp(2.0, 0.5, F)
adjacent_window(F) = round(lerp(8, 1, F))

since_last_s ← if last_retrieval_ts == 0 then +∞ else (now_ms() − last_retrieval_ts) / 1000
min_gap_s ← adjacent_window(F) × dt_ema
adjacent_ok ← (since_last_s ≥ min_gap_s)
force_check ← (drift_acc_pacing > max_wait_drift(F))

# Retrieval triggered when drift exceeds threshold, at memory boundary, or when drift exceeds max_wait_drift.
# Adjacent-window throttling is bypassed on boundaries and force_check.
if (drift_acc_pacing > pacing_thresh(S) OR should_flush OR force_check) AND
   (adjacent_ok OR should_flush OR force_check):
    trigger_check(); x_last_check ← x_t; drift_acc_pacing ← 0; last_retrieval_ts ← now_ms()
\end{verbatim}

High Sensitivity produces frequent checks triggered by small content
shifts; high Focus enforces strict drift limits. Memory boundaries
(Section~\ref{sec-boundary}) also trigger retrieval checks to ensure
context updates align with natural thought transitions.

\section{Experimental Results}\label{sec-experimental}

We present preliminary experimental results collected from live chat
sessions to validate the adaptive mechanisms.

\subsection{Threshold Adaptation}\label{threshold-adaptation}

The dynamic threshold (θ\_dynamic) successfully tracked score
distributions. In high-volatility inputs (within-accumulator drift\_acc
\textgreater{} 1.0), thresholds relaxed to \textasciitilde0.15, while
stable contexts tightened to \textasciitilde0.27.

\subsection{Boundary Detection}\label{boundary-detection}

Accumulator drift (drift\_acc) aligned with semantic shifts.
Conversation turns with distinct topics triggered flushes
(boundary\_score \textgreater{} 0.3) while coherent continuations
remained accumulated.

\subsection{Latency and Performance}\label{latency-and-performance}

End-to-end processing per token averaged \textless{} 50ms. Graph
expansion added \textless{} 10ms overhead due to efficient kNN (k=32)
and limited expansion depth (d=2).

\section{Implementation Considerations}\label{sec-implementation}

\subsection{Computational Complexity}\label{computational-complexity}

The primary per-stimulus cost is the k-nearest-neighbors (kNN) search
for metrics and retrieval. With \texttt{N} stored memories, naive search
is \texttt{O(N)}. We recommend maintaining an HNSW index (Malkov and
Yashunin 2018) to achieve \texttt{O(log\ N)} retrieval.

Recursive graph expansion adds overhead but is bounded by \texttt{depth}
and \texttt{branching\_factor}, effectively constant-time relative to
\texttt{N}. Total complexity per stimulus is \texttt{O(log\ N)} with
indexing.

\subsection{Memory Footprint}\label{memory-footprint}

The system stores: * \texttt{signal\_stream}: \texttt{n\_ctx} embeddings
(256d float32 = 1KB each). * \texttt{score\_stream}: Scalar history. *
\texttt{memory\_stream}: \texttt{N} stored memories (embedding +
metadata). * \texttt{graph}: Adjacency lists for graph edges.

For long-running instances, \texttt{memory\_stream} and \texttt{graph}
dominate. Pruning low-strength memories
(Section~\ref{sec-memory-strength}) is essential for bounding growth.

\subsection{Numerical Stability}\label{numerical-stability}

Exponential decays and probability computations can suffer from
underflow. We use \texttt{log-space} arithmetic where appropriate and
clamp all potentially divergent values (e.g., division by small sums)
with ε = 10⁻⁶.

\section{Conclusion and Future
Directions}\label{conclusion-and-future-directions}

We have presented Cortext, a three-knob adaptive memory architecture
that unifies working, episodic, and semantic memory processes under a
continuous control regime. By deriving most system parameters from
Focus, Sensitivity, and Stability, while fixing a small set of
invariants for controller stability, Cortext avoids the fragility of
hard-coded constants and enables robust operation across diverse
environments.

Key contributions include the formalization of knob-derived parameter
spaces, the integration of uncertainty-weighted Bayesian adaptation, and
the implementation of bio-inspired mechanisms like homeostatic rate
control and emotional consolidation.

Future work will focus on: 1. \textbf{Large-scale evaluation:}
Validating the architecture on long-horizon datasets (\textasciitilde1M+
steps). 2. \textbf{Multi-agent dynamics:} Exploring how Cortext
instances interact in collaborative settings. 3. \textbf{Hardware
acceleration:} Optimizing the kNN and graph traversal kernels for edge
devices.

Cortext represents a step toward more organic, life-long learning
systems that adapt continuously to their experiences, moving beyond
static knowledge bases toward truly cognitive memory.

\section*{References}\label{references}
\addcontentsline{toc}{section}{References}

\phantomsection\label{refs}
\begin{CSLReferences}{1}{0}
\bibitem[\citeproctext]{ref-anderson1994remembering}
Anderson, Michael C, Robert A Bjork, and Elizabeth L Bjork. 1994.
{``Remembering Can Cause Forgetting: Retrieval Dynamics in Long-Term
Memory.''} \emph{Journal of Experimental Psychology: Learning, Memory,
and Cognition} 20 (5): 1063.

\bibitem[\citeproctext]{ref-astrom2008feedback}
Åström, Karl Johan, and Richard M Murray. 2008. \emph{Feedback Systems:
An Introduction for Scientists and Engineers}. Princeton university
press.

\bibitem[\citeproctext]{ref-baddeley2000episodic}
Baddeley, Alan. 2000. {``The Episodic Buffer: A New Component of Working
Memory?''} \emph{Trends in Cognitive Sciences} 4 (11): 417--23.

\bibitem[\citeproctext]{ref-cowan2001magical}
Cowan, Nelson. 2001. {``The Magical Number 4 in Short-Term Memory: A
Reconsideration of Mental Storage Capacity.''} \emph{Behavioral and
Brain Sciences} 24 (1): 87--114.

\bibitem[\citeproctext]{ref-cowan2010magical}
---------. 2010. {``The Magical Mystery Four: How Is Working Memory
Capacity Limited, and Why?''} \emph{Current Directions in Psychological
Science} 19 (1): 51--57.

\bibitem[\citeproctext]{ref-hart1965memory}
Hart, J T. 1965. {``Memory and the Feeling-of-Knowing Experience.''}
\emph{Journal of Educational Psychology} 56 (4): 208.

\bibitem[\citeproctext]{ref-hunt1995subtlety}
Hunt, R Reed. 1995. {``The Subtlety of Distinctiveness: What von
Restorff Really Did.''} \emph{Psychonomic Bulletin \& Review} 2 (1):
105--12.

\bibitem[\citeproctext]{ref-labar2006cognitive}
LaBar, Kevin S, and Roberto Cabeza. 2006. {``Cognitive Neuroscience of
Emotional Memory.''} \emph{Nature Reviews Neuroscience} 7 (1): 54--64.

\bibitem[\citeproctext]{ref-liu1998sequential}
Liu, Jun S, and Rong Chen. 1998. {``Sequential Monte Carlo Methods for
Dynamic Systems.''} \emph{Journal of the American Statistical
Association} 93 (443): 1032--44.

\bibitem[\citeproctext]{ref-malkov2018efficient}
Malkov, Yu A, and Dmitry A Yashunin. 2018. {``Efficient and Robust
Approximate Nearest Neighbor Search Using Hierarchical Navigable Small
World Graphs.''} \emph{IEEE Transactions on Pattern Analysis and Machine
Intelligence} 42 (4): 824--36.

\bibitem[\citeproctext]{ref-mcclelland1995why}
McClelland, James L, Bruce L McNaughton, and Randall C O'Reilly. 1995.
{``Why There Are Complementary Learning Systems in the Hippocampus and
Neocortex: Insights from the Successes and Failures of Connectionist
Models of Learning and Memory.''} \emph{Psychological Review} 102 (3):
419.

\bibitem[\citeproctext]{ref-mccloskey1989catastrophic}
McCloskey, Michael, and Neal J Cohen. 1989. \emph{Catastrophic
Interference in Connectionist Networks: The Sequential Learning
Problem}. Vol. 24. Elsevier.

\bibitem[\citeproctext]{ref-mcgaugh2004amygdala}
McGaugh, James L. 2004. {``The Amygdala Modulates the Consolidation of
Memories of Emotionally Arousing Experiences.''} \emph{Annual Review of
Neuroscience} 27: 1--28.

\bibitem[\citeproctext]{ref-miller1956magical}
Miller, George A. 1956. {``The Magical Number Seven, Plus or Minus Two:
Some Limits on Our Capacity for Processing Information.''}
\emph{Psychological Review} 63 (2): 81.

\bibitem[\citeproctext]{ref-murdock1962serial}
Murdock Jr, Bennet B. 1962. {``The Serial Position Effect of Free
Recall.''} \emph{Journal of Experimental Psychology} 64 (5): 482.

\bibitem[\citeproctext]{ref-nader2003memory}
Nader, Karim. 2003. {``Memory Traces Unbound.''} \emph{Trends in
Neurosciences} 26 (2): 65--72.

\bibitem[\citeproctext]{ref-nader2000fear}
Nader, Karim, Glenn E Schafe, and Joseph E Le Doux. 2000. {``Fear
Memories Require Protein Synthesis in the Amygdala for Reconsolidation
After Retrieval.''} \emph{Nature} 406 (6797): 722--26.

\bibitem[\citeproctext]{ref-russell1980circumplex}
Russell, James A. 1980. {``A Circumplex Model of Affect.''}
\emph{Journal of Personality and Social Psychology} 39 (6): 1161.

\bibitem[\citeproctext]{ref-tulving1972episodic}
Tulving, Endel. 1972. {``Episodic and Semantic Memory.''}
\emph{Organization of Memory} 1: 381--403.

\end{CSLReferences}

\section*{Appendix A. State Variables Map}\label{sec-app-state}
\addcontentsline{toc}{section}{Appendix A. State Variables Map}

This appendix enumerates the state variables used by the specification
and separates retained state (carried across timesteps) from per-step
derived quantities. The canonical source for section references is
\texttt{docs/paper/sections/*.qmd}; anchors in these files are
authoritative when resolving cross-reference drift.

\subsection*{State + Initialization (Canonical)}\label{sec-state-init}
\addcontentsline{toc}{subsection}{State + Initialization (Canonical)}

On cold start (no persisted state), initialize retained state as follows
(unless otherwise specified by knob priors):

\begin{itemize}
\item
  \textbf{Global defaults:} \texttt{signals\_processed\ =\ 0},
  \texttt{u\_uncertainty\ =\ 0}, \texttt{mood\_vector\ =\ 0\_vector},
  \texttt{last\_mood\_ts\ =\ now\_ms()},
  \texttt{theta\_dynamic\ =\ theta\_target\ =\ θ\_prior(F,S,T)},
  \texttt{hysteresis\ =\ lerp(0.02,\ 0.25,\ T)},
  \texttt{half\_life\ =\ base\_half\_life(T)}, \texttt{m\_rate\ =\ 0},
  \texttt{rho\_hat\_prev\ =\ 0}, \texttt{dt\_ema\ =\ 0},
  \texttt{rate\_ticks\ =\ 0}, \texttt{reliability\ =\ 1},
  \texttt{retention\_ema\ =\ 0},
  \texttt{last\_rate\_timestamp\ =\ now\_ms()},
  \texttt{last\_retrieval\_ts\ =\ 0},
  \texttt{last\_embedding\ =\ unset}, \texttt{x\_pred\_ema\ =\ unset}.
\item
  \textbf{Accumulator defaults (per stream):}
  \texttt{μ\_acc\ =\ 0\_vector}, \texttt{drift\_acc\ =\ 0},
  \texttt{s\_sum\ =\ 0}, \texttt{s\_max\ =\ 0}, \texttt{n\ =\ 0},
  \texttt{e\_peak\ =\ 0\_vector}, \texttt{emo\_max\ =\ 0},
  \texttt{arousal\_sum\ =\ 0}, \texttt{acc\_signals\_window\ =\ {[}{]}},
  \texttt{t\_start\ =\ 0}, \texttt{last\_signal\_ts\ =\ 0},
  \texttt{last\_write\_ts\ =\ 0}, \texttt{eta\_acc\ =\ 0},
  \texttt{coherence\_prev\ =\ 0}, \texttt{drift\_accum\ =\ 0},
  \texttt{drift\_at\_last\_interrupt\ =\ 0},
  \texttt{drift\_acc\_pacing\ =\ 0}, \texttt{x\_last\_check\ =\ unset},
  \texttt{prev\_x\ =\ unset}.
\item
  \textbf{Per-memory defaults (on insert):} \texttt{strength\ =\ 1},
  \texttt{use\_frequency\ =\ 0}, \texttt{stability\ =\ 1},
  \texttt{connectivity\ =\ 0}, \texttt{drift\_mag\ =\ 0},
  \texttt{influence\ =\ 0}, \texttt{sustained\_influence\ =\ 0},
  \texttt{contextual\_gain\ =\ 0}, \texttt{redundancy\ =\ 0},
  \texttt{pre\_activation\ =\ 0}, \texttt{lability\_state\ =\ 0},
  \texttt{suppression\_count\ =\ 0}, \texttt{suppression\ =\ 0},
  \texttt{flashbulb\ =\ 0}, \texttt{s\_emotion\_max\ =\ 0},
  \texttt{s\_arousal\_avg\ =\ 0}.
\item
  \textbf{RLS defaults:} \texttt{w\_*\ =\ w\_bootstrap},
  \texttt{P\ =\ diag(1000)}, \texttt{blender\_ready\ =\ false},
  \texttt{blender\_update\_count\ =\ 0}.
\item
  \textbf{Knob-derived control parameters (initialized from priors):}
  \texttt{weight\_relevance}, \texttt{attention\_width},
  \texttt{coverage\_gain\_floor}, \texttt{mismatch\_weight},
  \texttt{weight\_novelty}, \texttt{weight\_surprise},
  \texttt{weight\_valence}, \texttt{weight\_arousal},
  \texttt{emotion\_gain}, \texttt{score\_gain}, \texttt{rate\_target},
  \texttt{rate\_decay}, \texttt{periphery\_half\_life},
  \texttt{salience\_half\_life}, \texttt{drift\_weight}.
\item
  \textbf{Buffers:} \texttt{signal\_stream}, \texttt{score\_stream},
  \texttt{memory\_stream}, \texttt{recent\_memory\_centroids} start
  empty.
\item
  \textbf{Retention history:} \texttt{retention\_history} starts empty
  (populated after the first signal step).
\item
  \textbf{Accumulator state (per stream; retained across timesteps):}

  \begin{itemize}
  \tightlist
  \item
    \texttt{\{μ\_acc,\ drift\_acc,\ s\_sum,\ s\_max,\ n,\ e\_peak,\ emo\_max,\ arousal\_sum,\ eta\_acc,\ coherence\_prev,\ acc\_signals\_window,\ t\_start,\ last\_signal\_ts,\ last\_write\_ts,\ drift\_accum,\ drift\_at\_last\_interrupt,\ drift\_acc\_pacing,\ x\_last\_check,\ prev\_x\}}
  \end{itemize}
\item
  \textbf{Global state (retained across timesteps):}

  \begin{itemize}
  \tightlist
  \item
    \texttt{\{signals\_processed,\ u\_uncertainty,\ mood\_vector,\ last\_mood\_ts,\ theta\_dynamic,\ theta\_target,\ hysteresis,\ half\_life,\ m\_rate,\ rho\_hat\_prev,\ dt\_ema,\ rate\_ticks,\ reliability,\ retention\_ema,\ last\_rate\_timestamp,\ last\_retrieval\_ts,\ last\_embedding,\ x\_pred\_ema,\ weight\_relevance,\ attention\_width,\ coverage\_gain\_floor,\ mismatch\_weight,\ weight\_novelty,\ weight\_surprise,\ weight\_valence,\ weight\_arousal,\ emotion\_gain,\ score\_gain,\ rate\_target,\ rate\_decay,\ periphery\_half\_life,\ salience\_half\_life,\ drift\_weight,\ blender\_ready,\ blender\_update\_count,\ blender\_P\}}
  \end{itemize}
\item
  \textbf{Buffers (retained across timesteps; bounded by window rules):}

  \begin{itemize}
  \tightlist
  \item
    \texttt{\{signal\_stream,\ score\_stream,\ memory\_stream,\ recent\_memory\_centroids,\ retention\_history\}}
  \end{itemize}
\item
  \textbf{Recorded signal fields (per signal):}

  \begin{itemize}
  \tightlist
  \item
    \texttt{\{coherence\_struct\_t\ →\ SIGNALS.coherence,\ focus\_spread\_t\ →\ SIGNALS.focus\_spread\}}
  \end{itemize}
\item
  \textbf{Recorded global fields:}

  \begin{itemize}
  \tightlist
  \item
    \texttt{\{u(t)\ →\ STATE.u\_uncertainty,\ M\_t\ →\ STATE.mood\_vector\}}
  \end{itemize}
\item
  \textbf{Per-step derived scalars/vectors (ephemeral; recomputed each
  step):}

  \begin{itemize}
  \tightlist
  \item
    \texttt{\{signal\_gap\_s,\ coherence\_curr,\ s\_avg,\ S\_window,\ boundary\_score,\ should\_flush,\ write\_memory,\ Δwrites,\ q\_retrieval\}}
  \end{itemize}
\end{itemize}

\section*{Appendix B. Derived Signals: Definitions and
Bounds}\label{sec-app-derived}
\addcontentsline{toc}{section}{Appendix B. Derived Signals: Definitions
and Bounds}

This appendix defines non-trivial derived signals used throughout the
specification. Unless otherwise stated, all derived scalars are clamped
to {[}0, 1{]}.

\textbf{novelty:} A normalized dissimilarity-to-context signal.

\begin{verbatim}
max_cos ← max_{c ∈ recent_context} cos(x_t, c)  # in [−1, 1]
novelty_t ← clamp((1 − max_cos) / 2, 0, 1)
\end{verbatim}

Fallback: if recent\_context is empty, set novelty\_t ← 1.

\textbf{μ\_sim:} A normalized mean similarity to context.

\begin{verbatim}
mean_cos ← mean_{c ∈ recent_context} cos(x_t, c)
μ_sim ← clamp((mean_cos + 1) / 2, 0, 1)
\end{verbatim}

Fallback: if recent\_context is empty, set μ\_sim ← 0.5.

\textbf{rarity\_t:} A normalized rarity signal defined as
dissimilarity-to-context mean.

\begin{verbatim}
rarity_t ← clamp(1 − μ_sim, 0, 1)
\end{verbatim}

Fallback: inherits μ\_sim fallback, so rarity\_t defaults to 0.5 when
recent\_context is empty.

\textbf{relevance\_to\_task(q, task\_ctx):} A normalized relevance of
embedding q to a task context set.

\begin{verbatim}
if |task_ctx| == 0: return 0.5
return clamp((cos(q, mean(task_ctx)) + 1) / 2, 0, 1)
\end{verbatim}

\textbf{novelty\_to\_set(q, S\_set\_embeddings):} A normalized novelty
of q relative to a set of embeddings.

\begin{verbatim}
novelty_to_set(q, S_set_embeddings) ← 1 − redundancy(q, S_set_embeddings)
\end{verbatim}

\textbf{ΔSSE:} A normalized improvement in reconstruction/prediction
error (utility proxy).

\begin{verbatim}
ΔSSE ← clamp((SSE_prev − SSE_curr) / max(SSE_prev, ε), 0, 1)
\end{verbatim}

Definitions:

\begin{verbatim}
SSE_curr ← ‖x_t − x_pred_t‖^2          # prediction error at t (x_pred_t = x_pred_ema for EMA predictor)
SSE_prev ← previous SSE_curr (t−1)     # or EWMA if smoothing is used
\end{verbatim}

Fallback: if no prediction model is available, or if dimensions
mismatch, set ΔSSE ← 0.

\textbf{redundancy(a, S\_set):} A normalized redundancy of item a w.r.t.
a set S\_set.

\begin{verbatim}
redundancy(a, S_set) ← max_{s ∈ S_set} clamp((cos(a, s) + 1) / 2, 0, 1)
\end{verbatim}

Fallback: if S\_set is empty, redundancy(a, S\_set) ← 0.

\textbf{coverage\_gain(candidate \textbar{} included\_set):} Incremental
coverage contribution of adding candidate.

\begin{verbatim}
coverage_gain(candidate | included_set) ← 1 − redundancy(candidate, included_set)
\end{verbatim}

\section*{Appendix C. Edge Case Rules}\label{sec-app-edge}
\addcontentsline{toc}{section}{Appendix C. Edge Case Rules}

These rules resolve empty-store and small-store cases to preserve
causality and avoid undefined metrics.

\begin{itemize}
\tightlist
\item
  \textbf{memory\_stream empty:} kNN returns empty; retrieval returns no
  candidates; graph traversal is skipped. Set
  \texttt{focus\_spread\_t\ =\ 0}.
\item
  \textbf{memory\_stream small:} use
  \texttt{k\_eff\ =\ min(k\_neighbors(T),\ \textbar{}memory\_stream\textbar{})}.
  If \texttt{k\_eff\ \textless{}\ 2}, set
  \texttt{focus\_spread\_t\ =\ 0}.
\item
  \textbf{graph absent or empty:} skip graph traversal; use vector
  search only.
\item
  \textbf{recent\_context empty:} apply fallbacks defined in Appendix B
  (e.g., \texttt{novelty\_t\ =\ 1}, \texttt{μ\_sim\ =\ 0.5},
  \texttt{rarity\_t\ =\ 0.5}, \texttt{μ\_ctx\ =\ 0\_vector} so
  \texttt{map01(cos)=0.5}).
\item
  \textbf{recent\_scores empty:} set \texttt{observed\_p90\ ←\ θ\_prior}
  in threshold updates.
\end{itemize}

\section*{Appendix D. Main Loop Execution Order and Normative
Invariants}\label{sec-app-loop}
\addcontentsline{toc}{section}{Appendix D. Main Loop Execution Order and
Normative Invariants}

This ordering is canonical and supersedes any other sequence described
elsewhere in the document.

Canonical single-step pseudocode (timestep t):

\begin{verbatim}
# Main loop (single timestep t)
now_s, now_ms, x_t ← read_inputs()
signal_stream.append(x_t)
recent_context ← tail(signal_stream, n_ctx(T))

# Structural metrics + uncertainty
coherence_struct_t, focus_spread_t, drift_mag_t, surprisal_t ← compute_structural_metrics(x_t)
x_pred_ema ← update_prediction_ema(x_pred_ema, x_t, T)  # after surprisal_t
u_t ← update_uncertainty(...)

# Adaptation + scoring
update_control_parameters(...)
score_t ← compute_composite_score(...)
score_stream.append(score_t)

# Threshold updates
θ_target ← prior_evidence_blend(...)
Δθ_* ← compute_threshold_deltas(...)
θ_dynamic ← update_theta_dynamic(...)
hysteresis ← update_hysteresis(...)

# Accumulator + boundary
update_accumulator(...)
boundary_score ← compute_boundary_score(...)
should_flush ← boundary_decision(boundary_score, time_caps, spike_bypass)
θ_memory ← θ_dynamic × M_write_refrac
write_memory ← force_write OR (should_flush AND (S_window > θ_memory))
Δwrites ← 1 if write_memory else 0  # computed immediately after write decision

# Post-write and retrieval (q_retrieval captured before any reset)
if write_memory: commit_memory_unit(); last_write_ts ← now_ms()
q_retrieval ← μ_acc  # cache current-unit query
streaming_pacing_check()
graph_retrieval(q_retrieval)
update_rate_state(Δwrites)
if should_flush: reset_accumulator()  # regardless of write outcome

# Interrupt gate (consumes retrieved candidates)
interrupt_gate_check()
\end{verbatim}

The normative execution order for a single timestep t:

\begin{enumerate}
\def\labelenumi{\arabic{enumi}.}
\tightlist
\item
  \textbf{Read Inputs:} \texttt{x\_t}, \texttt{now\_s()},
  \texttt{system\_time\_ms}.
\item
  \textbf{Update Buffers:} Append \texttt{x\_t} to
  \texttt{signal\_stream}. Update \texttt{n\_ctx(T)}. Retrieve
  \texttt{recent\_context}.
\item
  \textbf{Compute Structural Metrics:} \texttt{coherence\_struct},
  \texttt{focus\_spread}, \texttt{drift}, \texttt{surprisal}, then
  update \texttt{x\_pred\_ema} for the next step.
\item
  \textbf{Update Uncertainty:} \texttt{u(t)}.
\item
  \textbf{Compute Adaptation Dynamics:}

  \begin{itemize}
  \tightlist
  \item
    Update \texttt{α\_F(t)}, \texttt{α\_T(t)}.
  \item
    Update \texttt{weight\_relevance}, \texttt{half\_life}.
  \item
    Update \texttt{emotion} and \texttt{mood} state.
  \end{itemize}
\item
  \textbf{Compute Composite Score:}

  \begin{itemize}
  \tightlist
  \item
    Compute all 12 metrics.
  \item
    Update RLS weights \texttt{W\_blend}.
  \item
    Compute \texttt{score\_t}. Append to \texttt{score\_stream}.
  \end{itemize}
\item
  \textbf{Update Thresholds:}

  \begin{itemize}
  \tightlist
  \item
    Compute \texttt{θ\_target} (prior/evidence).
  \item
    Compute \texttt{Δθ} terms (homeo, sens, prec, emo, mood).
  \item
    Update \texttt{θ\_dynamic}. Update \texttt{hysteresis}.
  \end{itemize}
\item
  \textbf{Execute Memory Accumulation:}

  \begin{itemize}
  \tightlist
  \item
    Update accumulator (\texttt{drift\_acc}, \texttt{s\_max}, etc.).
  \item
    Compute \texttt{boundary\_score} and \texttt{should\_flush}.
  \item
    Check \texttt{spike\_bypass}.
  \item
    Compute \texttt{S\_window} and \texttt{θ\_memory}.
  \item
    Decide \texttt{write\_memory}.
  \end{itemize}
\item
  \textbf{Post-Write Updates and Retrieval:}

  \begin{itemize}
  \tightlist
  \item
    If \texttt{write\_memory}: Write to \texttt{memory\_stream}. Update
    \texttt{last\_write\_ts}.
  \item
    Cache \texttt{q\_retrieval\ ←\ μ\_acc} before any accumulator reset.
  \item
    Run streaming pacing and retrieval using \texttt{q\_retrieval}.
  \item
    Update \texttt{rate\_state} (homeostatic controller).
  \item
    If \texttt{should\_flush} (regardless of write): Call
    \texttt{reset\_accumulator()}.
  \end{itemize}
\item
  \textbf{Run Interrupt Gate:} Check for streaming interrupt using
  retrieved candidates (already filtered by write‑exclusion and WM
  overlap during retrieval).
\end{enumerate}

Timing notes: * \texttt{now\_s()} is captured at step 1 and reused for
all Δt computations in this timestep. * Threshold updates in step 7 use
the score computed in step 6. * Rate updates in step 9 use Δwrites from
the current step's \texttt{write\_memory} decision.

\textbf{Normative Invariants:} * \textbf{Causality:} Step \texttt{k}
uses only values computed in steps \texttt{1} through \texttt{k} or
retained from \texttt{t-1}. * \textbf{Write Atomicity:} A ``write'' is
an atomic commit of a \texttt{should\_flush} unit. Partial units are
never written. * \textbf{Gap Consistency:} \texttt{now\_s()} is fixed at
step 1. All \texttt{Δt} calculations use this fixed value.




\end{document}
